% abc80rom.tex -- den kommenterade disassembleringen av ABC80:s BASIC-ROM,
% byggd ur mallen (inifran/mall/bygg.sh abc80). Bara listningen och tabellerna
% som hör till den: systemvariablerna, felkoderna, tokens och symbolerna.
% Bokens löptext är inte med.

\documentclass[10pt,oneside]{report}
\usepackage{abc80stil}
\hypersetup{pdftitle={ABC80 BASIC-ROM, kommenterad disassemblering}}

\begin{document}
\hypersetup{pageanchor=false}
\begin{titlepage}
\centering
\vspace*{50mm}
{\Huge\bfseries ABC80 BASIC-ROM\par}
\vspace{6mm}
{\LARGE Kommenterad disassemblering\par}
\vspace{14mm}
{\large Z80-kod, \$0000--\$3FFF (16 KB), version 11273\par}
\vfill
{\small Hör till boken \emph{ABC80 inifrån}.\par}
\end{titlepage}
\hypersetup{pageanchor=true}
\pagenumbering{roman}
\tableofcontents
\cleardoublepage\pagenumbering{arabic}

\chapter{Inledning}
\par\noindent ABC80:s BASIC-ROM, 16 kilobyte Z80-kod på \adr{0000}–\adrx{3FFF}, i versionen med kontrollsumman 11273. Koden i listningen är inte publicerad med dokumentet utan fylldes i ur ROM:en när dokumentet byggdes; namnen och kommentarerna kommer ur arbetet med boken \emph{ABC80 inifrån}, där ROM:ens uppbyggnad och arbetssätt beskrivs.\par
\par\noindent Varje rad visar adressen, byten, eventuellt ett namn och instruktionen, följd av en kommentar. En rad visar högst fyra byte; har instruktionen eller tabellen fler står det tre punkter. Talen skrivs hexadecimalt med \$ framför och decimalt utan. Namnen på rutiner och variabler är valda för genomgången; ROM:en själv innehåller inga. Längre förklaringar står som block före den rutin de gäller. Instruktionernas namn är Zilogs.\par
\par\noindent Namnen och adresserna är länkar i PDF-filen. Osäkra tolkningar är markerade med (?). Uppgifter från andra källor har källan inom parentes; källförteckningen finns i boken.\par

% kapitel/variabler.tex -- fryst ur tex80.py 2026-10-06 (omstruktureringen steg 5);
% rättas här för hand.

\chapter{Systemvariabler och externa adresser}\label{kap:variabler}\index{systemvariabler|textbf}
\begingroup\setlength{\parskip}{5pt plus 1pt}
\par\noindent Kapitlet är en uppslagstabell till kapitlen om koden. Där en rutin läser eller skriver en fast adress i RAM står adressens namn i texten, och här står vad den används till. Namnen är valda för den här boken, som rutinernas namn; ROM:en själv innehåller inga namn. De tre första raderna är inga variabler utan adresser i tilläggs-ROM:arna, som ROM:en anropar eller hoppar till om en sådan krets sitter i (se avsnittet Minneskarta, s.~\pageref{avs:hardvaran:minneskarta}, och Anropen till tilläggskretsarna, s.~\pageref{avs:rom:extjp}).\par
\par\noindent Systemvariablerna ligger på \adr{FDD0}–\adr{FE3F}, och direkt efter dem kommer radbufferten \namn{BUF1} (se avsnittet Minnet, s.~\pageref{avs:arkitektur:minnet}). IY pekar hela tiden på \adr{FE16}, och de variabler som ROM:en når genom IY har förskjutningen inom parentes, till exempel IY+\adrx{0F} för \namn{FLAGS}. Adresser som inte står i tabellen används inte av ROM:en med eget namn.\par
\par\noindent De första variablerna delas av två delar av ROM:en. ASCII-aritmetiken (\kw{ADD\$}, \kw{SUB\$} och de andra, kapitlet ASCII-aritmetiken, s.~\pageref{kap:rom:ascii}) räknar med två tal skrivna som text, a och b, och sparar där var de börjar och slutar. Flyttalsfunktionerna (\kw{SIN}, \kw{LOG} och de andra, kapitlet Flyttalen, s.~\pageref{kap:rom:flyttal}) använder samma byte som arbetsminne. Eftersom de två aldrig räknar samtidigt kan de dela platsen.\par
\TabNeed{8\baselineskip}
\begin{longtable}{@{}l l p{0.62\textwidth}@{}}
\textbf{Namn} & \textbf{Adress} & \textbf{Beskrivning}\\\hline\endhead
\T{v:V405C}\texttt{V405C} & \texttt{\$405C} & tilläggs-ROM: målet för \kw{RST \$08}\\
\T{v:V604B}\texttt{V604B} & \texttt{\$604B} & tilläggs-ROM (skivminnets DOS): anropas vid kallstart om det står en \kw{JP}-instruktion (\adrx{C3}) här\\
\T{v:V6051}\texttt{V6051} & \texttt{\$6051} & tilläggs-ROM: anropas för att skriva ut ett felmeddelande, om bit 0 i \namn{EXTFLG} är satt\\
\T{v:ASBBEG}\texttt{ASBBEG} & \texttt{\$FDD0} & ASCII-aritmetiken: var b börjar; flyttalsfunktionerna: en kopia av argumentet x på \adr{FDD0}–\adr{FDD4} (\namn{FLDX})\\
\T{v:ASBEND}\texttt{ASBEND} & \texttt{\$FDD2} & ASCII-aritmetiken: var b slutar (byten \adrx{83} efter sista siffran)\\
\T{v:ASABEG}\texttt{ASABEG} & \texttt{\$FDD4} & ASCII-aritmetiken: var a börjar\\
\T{v:ASAEND}\texttt{ASAEND} & \texttt{\$FDD6} & ASCII-aritmetiken: var a slutar (byten \adrx{83})\\
\T{v:SINQ}\texttt{SINQ} & \texttt{\$FDE4} & \kw{SIN} och \kw{COS}: vinkelns kvadrant\\
\T{v:ASARG1}\texttt{ASARG1} & \texttt{\$FDE7} & ASCII-aritmetiken: var kopian av a\$ står i bufferten på stacken\\
\T{v:ASEND}\texttt{ASEND} & \texttt{\$FDE9} & ASCII-aritmetiken: där argumenten började på stacken; flyttalsfunktionerna: adressen där resultatet ska skrivas\\
\T{v:ASDEC}\texttt{ASDEC} & \texttt{\$FDEB} & ASCII-aritmetiken: n, antalet decimaler\\
\T{v:ASCODE}\texttt{ASCODE} & \texttt{\$FDED} & ASCII-aritmetiken: DE, platsen i uttryckets kod, sparad medan räkningen pågår\\
\T{v:NMICNT}\texttt{NMICNT} & \texttt{\$FDEF} & byten före \namnx{CLOCK}. \namn{NMICLK} läser BC från \adr{FDEF} för att få klockans låga byte i B och skriver tillbaka C oförändrad\\
\T{v:CLOCKH}\texttt{CLOCKH} & \texttt{\$FDF1} & \namnx{CLOCK}:s två höga byte (\adr{FDF1}–\adr{FDF2}), räknas ned var 256:e NMI; den låga byten på \adr{FDF0} räknas ned vid varje NMI\\
\T{v:CURROW}\texttt{CURROW} & \texttt{\$FDF3} & markörens rad (0–23)\\
\T{v:CURCOL}\texttt{CURCOL} & \texttt{\$FDF4} & markörens kolumn (0–39)\\
\T{v:KEYFLG}\texttt{KEYFLG} & \texttt{\$FDF5} & bit 7 = en tangent väntar (sätts av \namn{KEYINT})\\
\T{v:RPTDLY}\texttt{RPTDLY} & \texttt{\$FDF7} & autorepetitionen: tiden till nästa upprepning, i varv om 256; 70 efter en ny tangent, sedan 8\\
\T{v:CASST}\texttt{CASST} & \texttt{\$FDF8} & kassettläsningen: \adrx{FF} medan ett block tas emot, sedan resultatet: 0 = rätt, \adrx{A5} = formatfel (ETX saknas), \adrx{A3} = fel checksumma\\
\T{v:CASBUF}\texttt{CASBUF} & \texttt{\$FDF9} & kassetten: bufferten som läses till (0 = ingen läsning pågår)\\
\T{v:EDPTR}\texttt{EDPTR} & \texttt{\$FDFB} & \namn{INLINE}: mallraden som TAB-tangenten hämtar tecken ur\\
\T{v:CASBLK}\texttt{CASBLK} & \texttt{\$FDFD} & kassetten: blocknumret\\
\T{v:RCVPTR}\texttt{RCVPTR} & \texttt{\$FDFF} & kassettläsningen: DE, nästa byte i bufferten, sparad mellan avbrotten\\
\T{v:RCVBC}\texttt{RCVBC} & \texttt{\$FE01} & kassettläsningen: C = skiftregistret, B = antalet bitar kvar\\
\T{v:RCVCONT}\texttt{RCVCONT} & \texttt{\$FE03} & kassettläsningen: adressen där korutinen fortsätter vid nästa avbrott\\
\T{v:RCVSUM}\texttt{RCVSUM} & \texttt{\$FE05} & kassettläsningen: checksumman\\
\T{v:BRKFLG}\texttt{BRKFLG} & \texttt{\$FE07} & skilt från 0 om BREAK (CTRL-C) har tryckts\\
\T{v:DEVHEAD}\texttt{DEVHEAD} & \texttt{\$FE0A} & första posten i enhetslistan (avsnittet Enhetslistan, s.~\pageref{avs:arkitektur:enhetslistan})\\
\T{v:EXTFLG}\texttt{EXTFLG} & \texttt{\$FE0F} & bit 0 = tilläggs-ROM:en skriver felmeddelandena (via \adr{6051})\\
\T{v:BOFA}\texttt{BOFA} & \texttt{\$FE1C} & programmets början\\
\T{v:EOFA}\texttt{EOFA} & \texttt{\$FE1E} & programmets slut (slutmarkeringen \adrx{01})\\
\T{v:HEAP}\texttt{HEAP} & \texttt{\$FE20} & slutet av variablerna = första lediga byte på heapen\\
\T{v:FLAGS}\texttt{FLAGS} & \texttt{\$FE25} & flaggor: bit 4 = \kw{TRACE}, bit 5 = skyddat program (IY+\adrx{0F})\\
\T{v:STKTOP}\texttt{STKTOP} & \texttt{\$FE27} & stackens startvärde, som SP sätts till vid varje ny kommandorad och efter ett fel\\
\T{v:VARLST}\texttt{VARLST} & \texttt{\$FE29} & första posten i variabellistan (0 = inga variabler)\\
\T{v:VFE2B}\texttt{VFE2B} & \texttt{\$FE2B} & kortvalet: det senaste värdet till port 1, som skrivs till port 1 igen före varje \kw{OUT} och \kw{INP} (IY+\adrx{15})\\
\T{v:CURLIN}\texttt{CURLIN} & \texttt{\$FE2C} & adressen till raden som körs (eller hade fel)\\
\T{v:DATPTR}\texttt{DATPTR} & \texttt{\$FE2E} & \kw{READ}: nästa värde i \kw{DATA}-raderna\\
\T{v:DATLIN}\texttt{DATLIN} & \texttt{\$FE30} & \kw{READ}: \kw{DATA}-raden som läses\\
\T{v:FILLST}\texttt{FILLST} & \texttt{\$FE32} & fillistan: första FCB i kedjan (IY+\adrx{1C})\\
\T{v:VFE33}\texttt{VFE33} & \texttt{\$FE33} & \namn{FILLST}:s höga byte; 0 = inga öppna filer (IY+\adrx{1D})\\
\T{v:ONERR}\texttt{ONERR} & \texttt{\$FE34} & \kw{ON ERROR GOTO}: radens adress; höga byten 0 = avstängt (IY+\adrx{1E}/\adrx{1F})\\
\T{v:STMSP}\texttt{STMSP} & \texttt{\$FE36} & SP när satsen startade; \kw{ON ERROR} återställer stacken därifrån\\
\T{v:FNFRAM}\texttt{FNFRAM} & \texttt{\$FE3E} & ramen för det \kw{FN}-anrop som räknas, alltså var argumenten står\\
\T{v:BUF1}\texttt{BUF1} & \texttt{\$FE40} & radbufferten; medan en rad länkas (\namn{PRELNK}) pekare till det innersta öppna \kw{FOR}-blocket\\
\end{longtable}
\endgroup

% kapitel/listning.tex -- fryst ur tex80.py 2026-10-06 (omstruktureringen steg 5);
% rättas här för hand.

\chapter{Listning}
\input{listningar/abc80}

\appendix
% kapitel/felkoder.tex -- fryst ur tex80.py 2026-10-06 (omstruktureringen steg 5);
% rättas här för hand.

\chapter{Felkoder}\label{bil:felkoder}
\noindent När något går fel skriver ABC80 bara ett nummer, \texttt{ERR n}.
Tabellen ger för varje nummer texten och de rutiner i ROM:en där felet
kommer ifrån. Texterna\index{felkoder|textbf}\index{felhantering} står inte
i ROM:en utan kommer ur ABC80:s fellista.
\par\smallskip\noindent Felkoden är $\$80+n$. Den ges med \kw{RST \$10} följt av
koden, eller med ett hopp till en felingång, som är en \kw{RST \$10} med
koden efter. \emph{Fångas} betyder att felet står i listan
\namn{ERRTRP} och kan tas om hand av
\kw{ON ERROR GOTO}\index{ON ERROR GOTO@\kw{ON ERROR GOTO}} när ett program
körs. \namn{ERRTRP} har också 36, 39–41 och 43, som ROM:en själv aldrig
ger. \emph{Ingång} är felets ingång på \adrx{06F7}–\adrx{072B} (avsnittet
Felen, i boken). \emph{Ges i} är de namngivna rutiner
som ger felet, med en \kw{RST \$10} eller ett hopp till ingången; det kan
finnas fler ställen. Namnen länkar till raden i Rutinerna.\par\medskip
\TabNeed{30mm}
\begin{small}\setlength{\LTleft}{0pt}\setlength{\LTright}{\fill}\begin{longtable}{@{}r l >{\raggedright\arraybackslash\hyphenpenalty=10000}p{0.3\textwidth} l l >{\raggedright\arraybackslash}p{0.34\textwidth}@{}}
\textbf{Nr} & \textbf{Kod} & \textbf{Text} & \textbf{Fångas} & \textbf{Ingång} & \textbf{Ges i}\\\hline\endhead
0 & \texttt{\$80} & EJ TILLÅTET ÖKA {\char34}DIM{\char34} &  & \texttt{L06F9} & \hyperlink{l:DIM1}{\texttt{DIM1}}, \hyperlink{l:DIMSTR}{\texttt{DIMSTR}}\\
1 & \texttt{\$81} & FEL ANTAL INDEX &  & \texttt{L06FB} & \hyperlink{l:LNKSRC}{\texttt{LNKSRC}}\\
2 & \texttt{\$82} & OTILLÅTET SOM KOMMANDO &  & \texttt{L06FD} & \hyperlink{l:NODIR}{\texttt{NODIR}}, \hyperlink{l:LNKFOR}{\texttt{LNKFOR}}, \hyperlink{l:RETURN}{\texttt{RETURN}}\\
3 & \texttt{\$83} & MINNET FULLT &  & \texttt{L06FF} & \hyperlink{l:INSLIN}{\texttt{INSLIN}}, \hyperlink{l:CMDMERG}{\texttt{CMDMERG}}, \hyperlink{l:FN}{\texttt{FN}}, \hyperlink{l:CONCAT}{\texttt{CONCAT}}, \hyperlink{l:STRFIL}{\texttt{STRFIL}}, \hyperlink{l:EXVAL}{\texttt{EXVAL}}, \hyperlink{l:DIM1}{\texttt{DIM1}}, \hyperlink{l:DIM2}{\texttt{DIM2}}, \hyperlink{l:DIMSTR}{\texttt{DIMSTR}}, \hyperlink{l:MEMFUL}{\texttt{MEMFUL}}\\
4 & \texttt{\$84} & FÖR STORT FLYTTAL & ja & \texttt{L0701} & \hyperlink{l:POWI}{\texttt{POWI}}, \hyperlink{l:FADDC}{\texttt{FADDC}}, \hyperlink{l:FMUL}{\texttt{FMUL}}, \hyperlink{l:FMRND}{\texttt{FMRND}}, \hyperlink{l:FDIV}{\texttt{FDIV}}, \hyperlink{l:FDRND}{\texttt{FDRND}}, \hyperlink{l:FOVFL}{\texttt{FOVFL}}, \hyperlink{l:FEXP}{\texttt{FEXP}}, \hyperlink{l:FLOG}{\texttt{FLOG}}\\
5 & \texttt{\$85} & FÖR STORT INDEX &  & \texttt{L0703} & \hyperlink{l:VAR1}{\texttt{VAR1}}, \hyperlink{l:ERRIDX}{\texttt{ERRIDX}}\\
6 & \texttt{\$86} & HITTAR EJ DETTA RADNUMMER &  & \texttt{L0705} & \hyperlink{l:CMDED}{\texttt{CMDED}}, \hyperlink{l:REFADR}{\texttt{REFADR}}, \hyperlink{l:LNKELM}{\texttt{LNKELM}}\\
7 & \texttt{\$87} & FÖR STORT HELTAL & ja & \texttt{L0707} & \hyperlink{l:SMULDV}{\texttt{SMULDV}}, \hyperlink{l:FTIERR}{\texttt{FTIERR}}\\
8 & \texttt{\$88} & FINNS EJ I DETTA SYSTEM &  & \texttt{L0709} & \hyperlink{l:EXTJP1}{\texttt{EXTJP1}}, \hyperlink{l:EXTJP2}{\texttt{EXTJP2}}, \hyperlink{l:STEXE1}{\texttt{STEXE1}}, \hyperlink{l:STEXE2}{\texttt{STEXE2}}, \hyperlink{l:LNKELM}{\texttt{LNKELM}}, \hyperlink{l:STTAB1}{\texttt{STTAB1}}, \hyperlink{l:STTAB2}{\texttt{STTAB2}}, \hyperlink{l:FN}{\texttt{FN}}, \hyperlink{l:EXPTAB2}{\texttt{EXPTAB2}}\\
9 & \texttt{\$89} & INDEX UTANFÖR STRÄNGEN &  & \texttt{L070B} & \hyperlink{l:EXLEFT}{\texttt{EXLEFT}}, \hyperlink{l:STRSKP}{\texttt{STRSKP}}\\
10 & \texttt{\$8A} & TEXTEN FÅR EJ PLATS I STRÄNGEN &  & \texttt{L070D} & \hyperlink{l:ASGNS}{\texttt{ASGNS}}\\
11 & \texttt{\$8B} & FÖRSTÅR EJ &  &  & \hyperlink{l:TOKKW}{\texttt{TOKKW}}, \hyperlink{l:CMPPRT}{\texttt{CMPPRT}}, \hyperlink{l:VARNAM}{\texttt{VARNAM}}\\
12 & \texttt{\$8C} & FELAKTIGT TAL & ja & \texttt{L0713} & \hyperlink{l:RDVAL}{\texttt{RDVAL}}, \hyperlink{l:CMPNUM}{\texttt{CMPNUM}}, \hyperlink{l:EXVAL}{\texttt{EXVAL}}\\
13 & \texttt{\$8D} & FEL ANTAL ELLER TYP AV ARGUMENT &  &  & \hyperlink{l:LNKFN}{\texttt{LNKFN}}, \hyperlink{l:CMPDEF}{\texttt{CMPDEF}}, \hyperlink{l:CMPFN}{\texttt{CMPFN}}, \hyperlink{l:CMPINP}{\texttt{CMPINP}}, \hyperlink{l:CMPPOK}{\texttt{CMPPOK}}, \hyperlink{l:TYPCNV}{\texttt{TYPCNV}}, \hyperlink{l:CMPARG}{\texttt{CMPARG}}\\
14 & \texttt{\$8E} & OTILLÅTET TECKEN EFTER SATSEN &  &  & \hyperlink{l:TOKLIN}{\texttt{TOKLIN}}\\
15 & \texttt{\$8F} & {\char34}={\char34} SAKNAS ELLER PÅ FEL PLATS &  &  & \hyperlink{l:TOKKW}{\texttt{TOKKW}}, \hyperlink{l:CMPDEF}{\texttt{CMPDEF}}, \hyperlink{l:CMPLET}{\texttt{CMPLET}}\\
16 & \texttt{\$90} & RADNUMMER SAKNAS &  &  & \hyperlink{l:CMDMERG}{\texttt{CMDMERG}}, \hyperlink{l:CMPREF}{\texttt{CMPREF}}, \hyperlink{l:CMPON}{\texttt{CMPON}}\\
17 & \texttt{\$91} & OTILLÅTEN BLANDNING AV TAL OCH STRÄNGAR &  &  & \hyperlink{l:CMPSV}{\texttt{CMPSV}}, \hyperlink{l:CMPGET}{\texttt{CMPGET}}, \hyperlink{l:TOFLT}{\texttt{TOFLT}}, \hyperlink{l:CMPADD}{\texttt{CMPADD}}\\
18 & \texttt{\$92} & {\char34}){\char34} SAKNAS ELLER PÅ FEL PLATS &  &  & \hyperlink{l:CMPDEF}{\texttt{CMPDEF}}, \hyperlink{l:CMPFN}{\texttt{CMPFN}}, \hyperlink{l:PRIMARY}{\texttt{PRIMARY}}\\
19 & \texttt{\$93} & KAN EJ ÖPPNA FLER FILER & ja &  & \\
20 & \texttt{\$94} & FÖR LÅNG RAD &  &  & \hyperlink{l:TXTINP}{\texttt{TXTINP}}\\
21 & \texttt{\$95} & HITTAR EJ FILEN & ja & \texttt{L070F} & \hyperlink{l:OPEN}{\texttt{OPEN}}\\
22 & \texttt{\$96} & OTILLÅTEN SATS &  & \texttt{L0711} & \\
23 & \texttt{\$97} & {\char34}TO{\char34} SAKNAS &  &  & \hyperlink{l:CMPFOR}{\texttt{CMPFOR}}\\
24 & \texttt{\$98} & {\char34}NEXT{\char34} SAKNAS &  & \texttt{L0715} & \hyperlink{l:PRELNK}{\texttt{PRELNK}}\\
25 & \texttt{\$99} & FELAKTIG SATS EFTER {\char34}ON{\char34} &  &  & \hyperlink{l:CMPON}{\texttt{CMPON}}\\
26 & \texttt{\$9A} & FEL I ON-UTTRYCK &  & \texttt{L071D} & \hyperlink{l:ON}{\texttt{ON}}\\
27 & \texttt{\$9B} & {\char34}NEXT{\char34} UTAN {\char34}FOR{\char34} &  & \texttt{L0719} & \hyperlink{l:LNKNXT}{\texttt{LNKNXT}}\\
28 & \texttt{\$9C} & FEL VARIABEL EFTER {\char34}NEXT{\char34} &  & \texttt{L0717} & \hyperlink{l:LNKNXT}{\texttt{LNKNXT}}\\
29 & \texttt{\$9D} & {\char34}RETURN{\char34} UTAN {\char34}GOSUB{\char34} &  & \texttt{L071B} & \hyperlink{l:CMDRUN}{\texttt{CMDRUN}} (\adrx{0D6D})\\
30 & \texttt{\$9E} & DATA SLUT & ja & \texttt{L071F} & \hyperlink{l:READ}{\texttt{READ}}\\
31 & \texttt{\$9F} & FEL DATA TILL KOMMANDO &  & \texttt{L0721} & \hyperlink{l:CMDED}{\texttt{CMDED}}, \hyperlink{l:CMDLIST}{\texttt{CMDLIST}}, \hyperlink{l:CMDREN}{\texttt{CMDREN}}\\
32 & \texttt{\$A0} & FILEN EJ ÖPPNAD &  & \texttt{L0723} & \hyperlink{l:CURFIL}{\texttt{CURFIL}}\\
33 & \texttt{\$A1} & {\char34}AS FILE{\char34} SAKNAS &  &  & \hyperlink{l:CMPOPN}{\texttt{CMPOPN}}\\
34 & \texttt{\$A2} & SLUT PÅ FILEN & ja &  & \hyperlink{l:INPUT}{\texttt{INPUT}}\\
35 & \texttt{\$A3} & CHECKSUMMAFEL VID LÄSNING & ja &  & \hyperlink{l:RCVBLK}{\texttt{RCVBLK}}\\
37 & \texttt{\$A5} & FELAKTIGT RECORDFORMAT & ja &  & \hyperlink{l:CASOPN}{\texttt{CASOPN}}, \hyperlink{l:CASGET}{\texttt{CASGET}}, \hyperlink{l:RCVBLK}{\texttt{RCVBLK}}\\
38 & \texttt{\$A6} & RECORDNUMMER UTANFÖR FILEN & ja &  & \\
42 & \texttt{\$AA} & SKIVAN EJ KLAR & ja &  & \hyperlink{l:CASRBL}{\texttt{CASRBL}}\\
50 & \texttt{\$B2} & KVADRATROT UR NEGATIVT TAL &  & \texttt{L0725} & \hyperlink{l:FSQR}{\texttt{FSQR}}\\
51 & \texttt{\$B3} & ENHETEN UPPTAGEN &  &  & \\
52 & \texttt{\$B4} & EJ TILL DENNA ENHET &  & \texttt{L0727} & \hyperlink{l:CASTAB}{\texttt{CASTAB}}, \hyperlink{l:CMDMERG}{\texttt{CMDMERG}}, \hyperlink{l:CMDSAVE}{\texttt{CMDSAVE}}\\
53 & \texttt{\$B5} & FELAKTIG RAD &  & \texttt{L06F7} & \hyperlink{l:STEXE1}{\texttt{STEXE1}}\\
57 & \texttt{\$B9} & FUNKTIONEN EJ DEFINIERAD &  & \texttt{L0729} & \hyperlink{l:LNKFN}{\texttt{LNKFN}}\\
58 & \texttt{\$BA} & OGILTIGT TECKEN INLÄST &  &  & \hyperlink{l:TXTINP}{\texttt{TXTINP}}\\
59 & \texttt{\$BB} & FEL PROGRAMFORMAT &  &  & \hyperlink{l:CMDMERG}{\texttt{CMDMERG}}\\
60 & \texttt{\$BC} & BIT ADRESS 16 BITAR &  & \texttt{L072B} & \\
61 & \texttt{\$BD} & KOMMA SAKNAS &  &  & \\
62 & \texttt{\$BE} & DOT-ADRESS UTANFÖR SKÄRMEN &  &  & \hyperlink{l:DOTARG}{\texttt{DOTARG}}\\
63 & \texttt{\$BF} & {\char34}AS{\char34} SAKNAS &  &  & \\
64 & \texttt{\$C0} & FELAKTIG {\char34}RENAME{\char34} &  &  & \\
65 & \texttt{\$C1} & SPILL I ASCII-ARITMETIK &  & \hyperlink{l:ERR65}{\texttt{ERR65}} & \hyperlink{l:ASCARG}{\texttt{ASCARG}}\\
66 & \texttt{\$C2} & STRÄNG EJ NUMERISK &  & \hyperlink{l:ERR66}{\texttt{ERR66}} & \hyperlink{l:ASCAR}{\texttt{ASCAR}}, \hyperlink{l:ASCERR}{\texttt{ASCERR}}, \hyperlink{l:ASCARG}{\texttt{ASCARG}}\\
\end{longtable}\end{small}
\noindent Fyra av felen kommer inte där man väntar sig dem:\par
\begin{itemize}\setlength{\itemsep}{0pt}
\item Fel 1 kommer bara när samma namn står med ett index och med två
  (\kw{A(1)} och \kw{A(1,1)}). \kw{A} och \kw{A(1)}, eller \kw{A\$} och
  \kw{A\$(1)}, är olika variabler och går bra (avsnittet Variablerna,
  i boken).
\item Fel 13 är det tänkta felet för \kw{DIM A} och \kw{DIM A(1,2,3)}, men
  stacken är fel efter felet i \namn{CMPDMI}, och ABC80 ger i stället
  ett fel som beror på vad som råkar ligga i minnet, oftast 4 (avsnittet
  \kw{DIM}, i boken). Den senare ROM:en 9913 ger fel 13
  (i boken).
\item Fel 61 kommer aldrig: kontrollen av kommat i \namn{CMPDOT} slår
  inte till (avsnittet Grafikpunkterna, i boken).
\item Fel 62 kommer bara från \kw{SETDOT} och \kw{CLRDOT}. \kw{DOT}
  utanför skärmen ger inget fel, utan $-$1 eller 0 efter en byte i ROM:en
  (i boken).
\end{itemize}

% kapitel/tokens.tex -- fryst ur tex80.py 2026-10-06 (omstruktureringen steg 5);
% rättas här för hand.

\chapter{Tokens och tabeller}
\section{Direktkommandon}
\noindent Ur \hyperlink{l:CMDKW}{\texttt{CMDKW}} och
\hyperlink{l:CMDTAB}{\texttt{CMDTAB}}. Koden är index i CMDTAB\index{direktkommando}.\par\medskip
\TabNeed{30mm}
\begin{small}\setlength{\LTleft}{0pt}\setlength{\LTright}{\fill}\begin{longtable}{@{}l l l@{}}
\textbf{Kommando} & \textbf{Kod} & \textbf{Rutin}\\\hline\endhead
\texttt{BYE} & \texttt{\$80} & \hyperlink{l:L017A}{\texttt{L017A}} {\footnotesize(tillägg \texttt{\$604E})}\\
\texttt{RUN} & \texttt{\$81} & \hyperlink{l:CMDRUN}{\texttt{CMDRUN}}\\
\texttt{LIST} & \texttt{\$82} & \hyperlink{l:CMDLIST}{\texttt{CMDLIST}}\\
\texttt{NEW} & \texttt{\$83} & \hyperlink{l:NEW}{\texttt{NEW}}\\
\texttt{SCR} & \texttt{\$83} & \hyperlink{l:NEW}{\texttt{NEW}}\\
\texttt{CLEAR} & \texttt{\$84} & \hyperlink{l:CLEAR}{\texttt{CLEAR}}\\
\texttt{LOAD} & \texttt{\$85} & \hyperlink{l:CMDLOAD}{\texttt{CMDLOAD}}\\
\texttt{MERGE} & \texttt{\$86} & \hyperlink{l:CMDMERG}{\texttt{CMDMERG}}\\
\texttt{REN} & \texttt{\$87} & \hyperlink{l:CMDREN}{\texttt{CMDREN}}\\
\texttt{ED} & \texttt{\$88} & \hyperlink{l:CMDED}{\texttt{CMDED}}\\
\texttt{SAVE} & \texttt{\$89} & \hyperlink{l:CMDSAVE}{\texttt{CMDSAVE}}\\
\texttt{UNSAVE} & \texttt{\$8A} & \hyperlink{l:L014C}{\texttt{L014C}} {\footnotesize(tillägg \texttt{\$605A})}\\
\end{longtable}\end{small}
\section{Satser}
\noindent Token\index{token} $\$80$--$\$96$ ur \hyperlink{l:STKW1}{\texttt{STKW1}}.
Kompileringsrutinen står i \hyperlink{l:STTAB1}{\texttt{STTAB1}} och
körrutinen i \hyperlink{l:STEXE1}{\texttt{STEXE1}}, båda med index
token${}-\$80$. Satser i tabell 2 skrivs som $\$86$ följt av koden i
\hyperlink{l:STKW2}{\texttt{STKW2}} och har rutinerna i
\hyperlink{l:STTAB2}{\texttt{STTAB2}} och
\hyperlink{l:STEXE2}{\texttt{STEXE2}}.\par\medskip
\TabNeed{30mm}
\begin{small}\setlength{\LTleft}{0pt}\setlength{\LTright}{\fill}\begin{longtable}{@{}l l l l@{}}
\textbf{Token} & \textbf{Ord} & \textbf{Kompilering} & \textbf{Körning}\\\hline\endhead
\texttt{\$80} & \texttt{GOTO} & \hyperlink{l:CMPREF}{\texttt{CMPREF}} & \hyperlink{l:GOTO}{\texttt{GOTO}}\\
\texttt{\$81} & (\kw{GOTO} efter \kw{THEN}/\kw{ELSE}) & \hyperlink{l:CMPREF}{\texttt{CMPREF}} & \hyperlink{l:GOTO}{\texttt{GOTO}}\\
\texttt{\$82} & \texttt{LET} & \hyperlink{l:CMPLET}{\texttt{CMPLET}} & \hyperlink{l:LET}{\texttt{LET}}\\
\texttt{\$83} & (\kw{LET} utan ordet) & \hyperlink{l:CMPLET}{\texttt{CMPLET}} & \hyperlink{l:LET}{\texttt{LET}}\\
\texttt{\$84} & \texttt{PRINT} & \hyperlink{l:CMPPRT}{\texttt{CMPPRT}} & \hyperlink{l:PRINTS}{\texttt{PRINTS}}\\
\texttt{\$85} & \texttt{;} & \hyperlink{l:CMPPRT}{\texttt{CMPPRT}} & \hyperlink{l:PRINTS}{\texttt{PRINTS}}\\
\texttt{\$86} & (tabell 2, kod följer) & \hyperlink{l:RESET}{\texttt{RESET}} & \hyperlink{l:EXEST2}{\texttt{EXEST2}}\\
\texttt{\$87} & (hopp framåt, \kw{IF}) & \hyperlink{l:RESET}{\texttt{RESET}} & \hyperlink{l:EXESKP}{\texttt{EXESKP}}\\
\texttt{\$88} & \texttt{IF} & \hyperlink{l:CMPIF}{\texttt{CMPIF}} & \hyperlink{l:IF}{\texttt{IF}}\\
\texttt{\$89} & \texttt{INPUT} & \hyperlink{l:CMPINP}{\texttt{CMPINP}} & \hyperlink{l:INPT}{\texttt{INPT}}\\
\texttt{\$8A} & \texttt{INPUTLINE} & \hyperlink{l:CMPINL}{\texttt{CMPINL}} & \hyperlink{l:INPLIN}{\texttt{INPLIN}}\\
\texttt{\$8B} & \texttt{FOR} & \hyperlink{l:CMPFOR}{\texttt{CMPFOR}} & \hyperlink{l:FOR}{\texttt{FOR}}\\
\texttt{\$8C} & \texttt{NEXT} & \hyperlink{l:CMPNXT}{\texttt{CMPNXT}} & \hyperlink{l:NEXT}{\texttt{NEXT}}\\
\texttt{\$8D} & (inget ord) & \hyperlink{l:L0709}{fel 8} & \hyperlink{l:L0709}{fel 8}\\
\texttt{\$8E} & \texttt{READ} & \hyperlink{l:CMPRD}{\texttt{CMPRD}} & \hyperlink{l:READ}{\texttt{READ}}\\
\texttt{\$8F} & \texttt{RESTORE} & \hyperlink{l:CMPRES}{\texttt{CMPRES}} & \hyperlink{l:RESTOR}{\texttt{RESTOR}}\\
\texttt{\$90} & \texttt{GOSUB} & \hyperlink{l:CMPREF}{\texttt{CMPREF}} & \hyperlink{l:GOSUB}{\texttt{GOSUB}}\\
\texttt{\$91} & \texttt{RETURN} & \hyperlink{l:NOCODE}{\texttt{NOCODE}} & \hyperlink{l:RETURN}{\texttt{RETURN}}\\
\texttt{\$92} & \texttt{ON} & \hyperlink{l:CMPON}{\texttt{CMPON}} & \hyperlink{l:ON}{\texttt{ON}}\\
\texttt{\$93} & \texttt{DATA} & \hyperlink{l:CMPREM}{\texttt{CMPREM}} & \hyperlink{l:DATA}{\texttt{DATA}}\\
\texttt{\$94} & \texttt{ONERRORGOTO} & \hyperlink{l:CMPREF}{\texttt{CMPREF}} & \hyperlink{l:ONERRG}{\texttt{ONERRG}}\\
\texttt{\$95} & (felaktig rad) & \hyperlink{l:RESET}{\texttt{RESET}} & \hyperlink{l:L06F7}{fel 53}\\
\texttt{\$96} & \texttt{DEFFN} & \hyperlink{l:CMPDEF}{\texttt{CMPDEF}} & \hyperlink{l:DEF}{\texttt{DEF}}\\
\end{longtable}\end{small}
\section{Satser, tabell 2}
\TabNeed{30mm}
\begin{small}\setlength{\LTleft}{0pt}\setlength{\LTright}{\fill}\begin{longtable}{@{}l l l l@{}}
\textbf{Kod} & \textbf{Ord} & \textbf{Kompilering} & \textbf{Körning}\\\hline\endhead
\texttt{\$86 \$80} & \texttt{DIM} & \hyperlink{l:CMPDIM}{\texttt{CMPDIM}} & \hyperlink{l:DIM}{\texttt{DIM}}\\
\texttt{\$86 \$81} & \texttt{POKE} & \hyperlink{l:CMPPOK}{\texttt{CMPPOK}} & \hyperlink{l:POKE}{\texttt{POKE}}\\
\texttt{\$86 \$82} & \texttt{OUT} & \hyperlink{l:CMPOUT}{\texttt{CMPOUT}} & \hyperlink{l:OUTST}{\texttt{OUTST}}\\
\texttt{\$86 \$83} & (inget ord) & \hyperlink{l:CMPINX}{\texttt{CMPINX}} & \hyperlink{l:L0709}{fel 8}\\
\texttt{\$86 \$84} & \texttt{REM} & \hyperlink{l:CMPREM}{\texttt{CMPREM}} & \hyperlink{l:REM}{\texttt{REM}}\\
\texttt{\$86 \$85} & \texttt{OPEN} & \hyperlink{l:CMPOPN}{\texttt{CMPOPN}} & \hyperlink{l:OPENST}{\texttt{OPENST}}\\
\texttt{\$86 \$86} & \texttt{PREPARE} & \hyperlink{l:CMPOPN}{\texttt{CMPOPN}} & \hyperlink{l:PREPST}{\texttt{PREPST}}\\
\texttt{\$86 \$87} & \texttt{CLOSE} & \hyperlink{l:CMPCLS}{\texttt{CMPCLS}} & \hyperlink{l:CLOSST}{\texttt{CLOSST}}\\
\texttt{\$86 \$88} & \texttt{RANDOMIZE} & \hyperlink{l:CMPNOP}{\texttt{CMPNOP}} & \hyperlink{l:RANDOM}{\texttt{RANDOM}}\\
\texttt{\$86 \$89} & \texttt{STOP} & \hyperlink{l:NOCODE}{\texttt{NOCODE}} & \hyperlink{l:STOP}{\texttt{STOP}}\\
\texttt{\$86 \$8A} & \texttt{END} & \hyperlink{l:NOCODE}{\texttt{NOCODE}} & \hyperlink{l:ENDPRG}{\texttt{ENDPRG}}\\
\texttt{\$86 \$8B} & \texttt{CMD} & \hyperlink{l:L0147}{\texttt{L0147}} {\footnotesize(tillägg \texttt{\$7000})} & \hyperlink{l:L0166}{\texttt{L0166}} {\footnotesize(tillägg \texttt{\$7003})}\\
\texttt{\$86 \$8C} & (inget ord) & \hyperlink{l:L0709}{fel 8} & \hyperlink{l:L0709}{fel 8}\\
\texttt{\$86 \$8D} & (inget ord) & \hyperlink{l:L0709}{fel 8} & \hyperlink{l:L0709}{fel 8}\\
\texttt{\$86 \$8E} & \texttt{SETDOT} & \hyperlink{l:CMPDOT}{\texttt{CMPDOT}} & \hyperlink{l:SETDOT}{\texttt{SETDOT}}\\
\texttt{\$86 \$8F} & \texttt{CLRDOT} & \hyperlink{l:CMPDOT}{\texttt{CMPDOT}} & \hyperlink{l:CLRDOT}{\texttt{CLRDOT}}\\
\texttt{\$86 \$90} & \texttt{GET} & \hyperlink{l:CMPGET}{\texttt{CMPGET}} & \hyperlink{l:GET}{\texttt{GET}}\\
\texttt{\$86 \$91} & \texttt{CHAIN} & \hyperlink{l:CMPSTX}{\texttt{CMPSTX}} & \hyperlink{l:CHAIN}{\texttt{CHAIN}}\\
\texttt{\$86 \$92} & \texttt{KILL} & \hyperlink{l:CMPSTX}{\texttt{CMPSTX}} & \hyperlink{l:L0170}{\texttt{L0170}} {\footnotesize(tillägg \texttt{\$605D})}\\
\texttt{\$86 \$93} & \texttt{NAME} & \hyperlink{l:L0151}{\texttt{L0151}} {\footnotesize(tillägg \texttt{\$6054})} & \hyperlink{l:L0175}{\texttt{L0175}} {\footnotesize(tillägg \texttt{\$6057})}\\
\texttt{\$86 \$94} & \texttt{TRACE} & \hyperlink{l:NOCODE}{\texttt{NOCODE}} & \hyperlink{l:TRACE}{\texttt{TRACE}}\\
\texttt{\$86 \$95} & \texttt{NOTRACE} & \hyperlink{l:NOCODE}{\texttt{NOCODE}} & \hyperlink{l:NOTRAC}{\texttt{NOTRAC}}\\
\end{longtable}\end{small}
\section{Övriga nyckelord}
\noindent Tabeller med ett ord\index{nyckelord} var, som kompileringsrutinerna söker i.
Token skrivs i satsens kod och listas genom STKW1.\par\medskip
\TabNeed{30mm}
\begin{small}\setlength{\LTleft}{0pt}\setlength{\LTright}{\fill}\begin{longtable}{@{}l l l@{}}
\textbf{Tabell} & \textbf{Token} & \textbf{Ord}\\\hline\endhead
\hyperlink{l:KWTHEN}{\texttt{KWTHEN}} & \texttt{\$BD} & \texttt{THEN}\\
\hyperlink{l:KWELSE}{\texttt{KWELSE}} & \texttt{\$BC} & \texttt{ELSE}\\
\hyperlink{l:KWSTEP}{\texttt{KWSTEP}} & \texttt{\$BB} & \texttt{STEP}\\
\hyperlink{l:KWTO}{\texttt{KWTO}} & \texttt{\$BA} & \texttt{TO}\\
\hyperlink{l:KWASF}{\texttt{KWASF}} & \texttt{\$B8} & \texttt{ASFILE}\\
\hyperlink{l:KWAS}{\texttt{KWAS}} & \texttt{\$B7} & \texttt{AS}\\
\end{longtable}\end{small}
\section{Operatorer}
\noindent Nivåerna i uttryckskompilatorn, från lägst till högst
prioritet\index{prioritet (operatorer)} (\hyperlink{l:PREC1}{\texttt{PREC1}}). Aritmetik och jämförelser
har en token per typ: bas${}+0$ flyttal, ${}+1$ heltal, ${}+2$ sträng;
körrutinerna står i \hyperlink{l:EXPTAB1}{\texttt{EXPTAB1}}.\par\medskip
\TabNeed{30mm}
\begin{small}\setlength{\LTleft}{0pt}\setlength{\LTright}{\fill}\begin{longtable}{@{}l l l l l >{\raggedright\arraybackslash}p{0.3\textwidth}@{}}
\textbf{Nivå} & \textbf{Operator} & \textbf{Token} & \textbf{Klass} & \textbf{Kompilering} & \textbf{Körning}\\\hline\endhead
\hyperlink{l:PREC1}{\texttt{PREC1}} & \texttt{EQV} & \texttt{\$DC} & logisk & \hyperlink{l:PRECLV}{\texttt{PRECLV}} & \hyperlink{l:OPEQV}{\texttt{OPEQV}}\\
\hyperlink{l:PREC2}{\texttt{PREC2}} & \texttt{IMP} & \texttt{\$DD} & logisk & \hyperlink{l:PRECLV}{\texttt{PRECLV}} & \hyperlink{l:OPIMP}{\texttt{OPIMP}}\\
\hyperlink{l:PREC3}{\texttt{PREC3}} & \texttt{OR} & \texttt{\$DE} & logisk & \hyperlink{l:PRECLV}{\texttt{PRECLV}} & \hyperlink{l:OPOR}{\texttt{OPOR}}\\
\hyperlink{l:PREC3}{\texttt{PREC3}} & \texttt{XOR} & \texttt{\$DF} & logisk & \hyperlink{l:PRECLV}{\texttt{PRECLV}} & \hyperlink{l:OPXOR}{\texttt{OPXOR}}\\
\hyperlink{l:PREC4}{\texttt{PREC4}} & \texttt{AND} & \texttt{\$E0} & logisk & \hyperlink{l:PRECLV}{\texttt{PRECLV}} & \hyperlink{l:OPAND}{\texttt{OPAND}}\\
\hyperlink{l:PREC5}{\texttt{PREC5}} & \texttt{NOT} & \texttt{\$E1} & unär & \hyperlink{l:CMPNOT}{\texttt{CMPNOT}} & \hyperlink{l:OPNOT}{\texttt{OPNOT}}\\
\hyperlink{l:PREC6}{\texttt{PREC6}} & \texttt{{\char60}=} & \texttt{\$F1} & jämförelse & \hyperlink{l:PRECLV}{\texttt{PRECLV}} & \hyperlink{l:LEF}{\texttt{LEF}} \hyperlink{l:LEI}{\texttt{LEI}} \hyperlink{l:LES}{\texttt{LES}}\\
\hyperlink{l:PREC6}{\texttt{PREC6}} & \texttt{{\char60}{\char62}} & \texttt{\$E5} & jämförelse & \hyperlink{l:PRECLV}{\texttt{PRECLV}} & \hyperlink{l:NEF}{\texttt{NEF}} \hyperlink{l:NEI}{\texttt{NEI}} \hyperlink{l:NES}{\texttt{NES}}\\
\hyperlink{l:PREC6}{\texttt{PREC6}} & \texttt{{\char60}} & \texttt{\$E8} & jämförelse & \hyperlink{l:PRECLV}{\texttt{PRECLV}} & \hyperlink{l:LTF}{\texttt{LTF}} \hyperlink{l:LTI}{\texttt{LTI}} \hyperlink{l:LTS}{\texttt{LTS}}\\
\hyperlink{l:PREC6}{\texttt{PREC6}} & \texttt{{\char62}=} & \texttt{\$EB} & jämförelse & \hyperlink{l:PRECLV}{\texttt{PRECLV}} & \hyperlink{l:GEF}{\texttt{GEF}} \hyperlink{l:GEI}{\texttt{GEI}} \hyperlink{l:GES}{\texttt{GES}}\\
\hyperlink{l:PREC6}{\texttt{PREC6}} & \texttt{{\char62}} & \texttt{\$EE} & jämförelse & \hyperlink{l:PRECLV}{\texttt{PRECLV}} & \hyperlink{l:GTF}{\texttt{GTF}} \hyperlink{l:GTI}{\texttt{GTI}} \hyperlink{l:GTS}{\texttt{GTS}}\\
\hyperlink{l:PREC6}{\texttt{PREC6}} & \texttt{=} & \texttt{\$E2} & jämförelse & \hyperlink{l:PRECLV}{\texttt{PRECLV}} & \hyperlink{l:EQF}{\texttt{EQF}} \hyperlink{l:EQI}{\texttt{EQI}} \hyperlink{l:EQS}{\texttt{EQS}}\\
\hyperlink{l:PREC7}{\texttt{PREC7}} & \texttt{+} & \texttt{\$F4} & aritmetik & \hyperlink{l:CMPADD}{\texttt{CMPADD}} & \hyperlink{l:ADDF}{\texttt{ADDF}} \hyperlink{l:ADDI}{\texttt{ADDI}} \hyperlink{l:CONCAT}{\texttt{CONCAT}}\\
\hyperlink{l:PREC7}{\texttt{PREC7}} & \texttt{-} & \texttt{\$F7} & aritmetik & \hyperlink{l:CMPADD}{\texttt{CMPADD}} & \hyperlink{l:SUBF}{\texttt{SUBF}} \hyperlink{l:SUBI}{\texttt{SUBI}}\\
\hyperlink{l:PREC8}{\texttt{PREC8}} & \texttt{*} & \texttt{\$F9} & aritmetik & \hyperlink{l:PRECLV}{\texttt{PRECLV}} & \hyperlink{l:MULF}{\texttt{MULF}} \hyperlink{l:MULI}{\texttt{MULI}}\\
\hyperlink{l:PREC8}{\texttt{PREC8}} & \texttt{/} & \texttt{\$FB} & aritmetik & \hyperlink{l:PRECLV}{\texttt{PRECLV}} & \hyperlink{l:DIVF}{\texttt{DIVF}} \hyperlink{l:DIVI}{\texttt{DIVI}}\\
\hyperlink{l:PREC9}{\texttt{PREC9}} & \texttt{\^{}}, \texttt{**} & \texttt{\$FD} & potens & \hyperlink{l:PRECLV}{\texttt{PRECLV}} & \hyperlink{l:POWF}{\texttt{POWF}} \hyperlink{l:POWI}{\texttt{POWI}} \hyperlink{l:POWFI}{\texttt{POWFI}}\\
\hyperlink{l:PREC10}{\texttt{PREC10}} &  &  & primärt uttryck & \hyperlink{l:PRIMARY}{\texttt{PRIMARY}} & \\
\end{longtable}\end{small}
\section{Funktioner}
\noindent Token $\$80+k$ ur \hyperlink{l:FNKW}{\texttt{FNKW}};
koden $k$ är index i \hyperlink{l:EXPTAB2}{\texttt{EXPTAB2}}. Formen och
typerna är ur \hyperlink{l:FNTYPE}{\texttt{FNTYPE}}: \texttt{x} flyttal,
\texttt{i} heltal, \texttt{s\$} sträng, \texttt{[ ]} får utelämnas.\par\medskip
\TabNeed{30mm}
\begin{small}\setlength{\LTleft}{0pt}\setlength{\LTright}{\fill}\begin{longtable}{@{}l l l l@{}}
\textbf{Kod} & \textbf{Form} & \textbf{Resultat} & \textbf{Rutin}\\\hline\endhead
\texttt{\$00} & \texttt{FN} & -- & \hyperlink{l:FN}{\texttt{FN}}\\
\texttt{\$01} & \texttt{ABS(x)} & flyttal & \hyperlink{l:EXABS}{\texttt{EXABS}}\\
\texttt{\$02} & \texttt{ATN(x)} & flyttal & \hyperlink{l:EXATN}{\texttt{EXATN}}\\
\texttt{\$03} & \texttt{COS(x)} & flyttal & \hyperlink{l:EXCOS}{\texttt{EXCOS}}\\
\texttt{\$04} & \texttt{EXP(x)} & flyttal & \hyperlink{l:EXEXP1}{\texttt{EXEXP1}}\\
\texttt{\$05} & \texttt{FIX(x)} & flyttal & \hyperlink{l:EXFIX}{\texttt{EXFIX}}\\
\texttt{\$06} & \texttt{INT(x)} & flyttal & \hyperlink{l:EXINT}{\texttt{EXINT}}\\
\texttt{\$07} & \texttt{LOG(x)} & flyttal & \hyperlink{l:EXLOG}{\texttt{EXLOG}}\\
\texttt{\$08} & \texttt{LOG10(x)} & flyttal & \hyperlink{l:EXLG10}{\texttt{EXLG10}}\\
\texttt{\$09} & \texttt{PI} & flyttal & \hyperlink{l:EXPI}{\texttt{EXPI}}\\
\texttt{\$0A} & \texttt{RND[(x)]} & flyttal & \hyperlink{l:EXRND}{\texttt{EXRND}}\\
\texttt{\$0B} & \texttt{SGN(x)} & heltal & \hyperlink{l:EXSGN}{\texttt{EXSGN}}\\
\texttt{\$0C} & \texttt{SIN(x)} & flyttal & \hyperlink{l:EXSIN}{\texttt{EXSIN}}\\
\texttt{\$0D} & \texttt{SQR(x)} & flyttal & \hyperlink{l:EXSQR}{\texttt{EXSQR}}\\
\texttt{\$0E} & \texttt{TAN(x)} & flyttal & \hyperlink{l:EXTAN}{\texttt{EXTAN}}\\
\texttt{\$0F} & \texttt{ASC(s\$)} & heltal & \hyperlink{l:EXASC}{\texttt{EXASC}}\\
\texttt{\$10} & \texttt{CHR\$(i[,i[,i[,i]]])} & sträng & \hyperlink{l:EXCHR}{\texttt{EXCHR}}\\
\texttt{\$11} & \texttt{LEFT\$(s\$,i)} & sträng & \hyperlink{l:EXLEFT}{\texttt{EXLEFT}}\\
\texttt{\$12} & \texttt{RIGHT\$(s\$,i)} & sträng & \hyperlink{l:EXRGHT}{\texttt{EXRGHT}}\\
\texttt{\$13} & \texttt{MID\$(s\$,i,i)} & sträng & \hyperlink{l:EXMID}{\texttt{EXMID}}\\
\texttt{\$14} & \texttt{LEN(s\$)} & heltal & \hyperlink{l:EXLEN}{\texttt{EXLEN}}\\
\texttt{\$15} & \texttt{INSTR(i,s\$,s\$)} & heltal & \hyperlink{l:EXINST}{\texttt{EXINST}}\\
\texttt{\$16} & \texttt{SPACE\$(i)} & sträng & \hyperlink{l:EXSPC}{\texttt{EXSPC}}\\
\texttt{\$17} & \texttt{STRING\$(i,i)} & sträng & \hyperlink{l:EXSTRG}{\texttt{EXSTRG}}\\
\texttt{\$18} & \texttt{NUM\$(x)} & sträng & \hyperlink{l:EXNUM}{\texttt{EXNUM}}\\
\texttt{\$19} & \texttt{VAL(s\$)} & flyttal & \hyperlink{l:EXVAL}{\texttt{EXVAL}}\\
\texttt{\$1A} & \texttt{SWAP\%(i)} & heltal & \hyperlink{l:EXSWAP}{\texttt{EXSWAP}}\\
\texttt{\$1B} & \texttt{PEEK(i)} & heltal & \hyperlink{l:EXPEEK}{\texttt{EXPEEK}}\\
\texttt{\$1C} & \texttt{INP(i)} & heltal & \hyperlink{l:EXINP}{\texttt{EXINP}}\\
\texttt{\$1D} & \texttt{CALL(i[,i])} & heltal & \hyperlink{l:EXCALL}{\texttt{EXCALL}}\\
\texttt{\$1E} & \texttt{ERRCODE} & heltal & \hyperlink{l:EXERR}{\texttt{EXERR}}\\
\texttt{\$1F} & \texttt{IEC\$(i)} & sträng & \hyperlink{l:EXIEC}{\texttt{EXIEC}} {\footnotesize(tillägg \texttt{\$7006})}\\
\texttt{\$20} & \texttt{TAB(i)} & sträng & \hyperlink{l:EXTAB}{\texttt{EXTAB}}\\
\texttt{\$21} & \texttt{CUR(i,i)} & sträng & \hyperlink{l:EXCUR}{\texttt{EXCUR}}\\
\texttt{\$23} & \texttt{DOT(i,i)} & heltal & \hyperlink{l:EXDOT}{\texttt{EXDOT}}\\
\texttt{\$30} & \texttt{ADD\$(s\$,s\$,i)} & sträng & \hyperlink{l:EXADD}{\texttt{EXADD}}\\
\texttt{\$31} & \texttt{SUB\$(s\$,s\$,i)} & sträng & \hyperlink{l:EXSUB}{\texttt{EXSUB}}\\
\texttt{\$32} & \texttt{MUL\$(s\$,s\$,i)} & sträng & \hyperlink{l:EXMUL}{\texttt{EXMUL}}\\
\texttt{\$33} & \texttt{DIV\$(s\$,s\$,i)} & sträng & \hyperlink{l:EXDIV}{\texttt{EXDIV}}\\
\texttt{\$34} & \texttt{COMP\%(s\$,s\$)} & heltal & \hyperlink{l:EXCOMP}{\texttt{EXCOMP}}\\
\end{longtable}\end{small}
\noindent Koder utan namn i FNKW: \texttt{\$22}, \texttt{\$24}, \texttt{\$25}, \texttt{\$26}, \texttt{\$27}, \texttt{\$28}, \texttt{\$29}, \texttt{\$2A}, \texttt{\$2B}, \texttt{\$2C}, \texttt{\$2D}, \texttt{\$2E}, \texttt{\$2F} (\hyperlink{l:L0709}{fel 8}).\par
\section{Uttryckskod}
\noindent Token $\$B6$--$\$FF$ i den kompilerade uttryckskoden och
rutinen i \hyperlink{l:EXPTAB1}{\texttt{EXPTAB1}} som utför den. Vad varje
token betyder beskrivs vid tabellen i listningen.\par\medskip
\begin{multicols}{3}\small\raggedright
\noindent\texttt{\$B6}\quad \hyperlink{l:ASGNF}{\texttt{ASGNF}}\par
\noindent\texttt{\$B7}\quad \hyperlink{l:ASGNI}{\texttt{ASGNI}}\par
\noindent\texttt{\$B8}\quad \hyperlink{l:ASGNS}{\texttt{ASGNS}}\par
\noindent\texttt{\$B9}\quad \hyperlink{l:ENDF}{\texttt{ENDF}}\par
\noindent\texttt{\$BA}\quad \hyperlink{l:EXEND}{\texttt{EXEND}}\par
\noindent\texttt{\$BB}\quad \hyperlink{l:ENDS}{\texttt{ENDS}}\par
\noindent\texttt{\$BC}\quad \hyperlink{l:ENDVF}{\texttt{ENDVF}}\par
\noindent\texttt{\$BD}\quad \hyperlink{l:EXEND}{\texttt{EXEND}}\par
\noindent\texttt{\$BE}\quad \hyperlink{l:ENDVS}{\texttt{ENDVS}}\par
\noindent\texttt{\$BF}\quad \hyperlink{l:ENDDIM}{\texttt{ENDDIM}}\par
\noindent\texttt{\$C0}\quad \hyperlink{l:VAR0}{\texttt{VAR0}}\par
\noindent\texttt{\$C1}\quad \hyperlink{l:VAR0}{\texttt{VAR0}}\par
\noindent\texttt{\$C2}\quad \hyperlink{l:VAR0}{\texttt{VAR0}}\par
\noindent\texttt{\$C3}\quad \hyperlink{l:CONSTF}{\texttt{CONSTF}}\par
\noindent\texttt{\$C4}\quad \hyperlink{l:VAR1F}{\texttt{VAR1F}}\par
\noindent\texttt{\$C5}\quad \hyperlink{l:VAR1I}{\texttt{VAR1I}}\par
\noindent\texttt{\$C6}\quad \hyperlink{l:VAR1S}{\texttt{VAR1S}}\par
\noindent\texttt{\$C7}\quad \hyperlink{l:CONSTI}{\texttt{CONSTI}}\par
\noindent\texttt{\$C8}\quad \hyperlink{l:VAR2F}{\texttt{VAR2F}}\par
\noindent\texttt{\$C9}\quad \hyperlink{l:VAR2I}{\texttt{VAR2I}}\par
\noindent\texttt{\$CA}\quad \hyperlink{l:VAR2S}{\texttt{VAR2S}}\par
\noindent\texttt{\$CB}\quad \hyperlink{l:CONSTS}{\texttt{CONSTS}}\par
\noindent\texttt{\$CC}\quad \hyperlink{l:DIMELM}{\texttt{DIMELM}}\par
\noindent\texttt{\$CD}\quad \hyperlink{l:EXFUNC}{\texttt{EXFUNC}}\par
\noindent\texttt{\$CE}\quad \hyperlink{l:EXFUNC}{\texttt{EXFUNC}}\par
\noindent\texttt{\$CF}\quad \hyperlink{l:EXFUNC}{\texttt{EXFUNC}}\par
\noindent\texttt{\$D0}\quad \hyperlink{l:EXFUNC}{\texttt{EXFUNC}}\par
\noindent\texttt{\$D1}\quad \hyperlink{l:EXFUNC}{\texttt{EXFUNC}}\par
\noindent\texttt{\$D2}\quad \hyperlink{l:PUSH0}{\texttt{PUSH0}}\par
\noindent\texttt{\$D3}\quad \hyperlink{l:EXEXP}{\texttt{EXEXP}}\par
\noindent\texttt{\$D4}\quad \hyperlink{l:GETF}{\texttt{GETF}}\par
\noindent\texttt{\$D5}\quad \hyperlink{l:GETI}{\texttt{GETI}}\par
\noindent\texttt{\$D6}\quad \hyperlink{l:GETS}{\texttt{GETS}}\par
\noindent\texttt{\$D7}\quad \hyperlink{l:NEGF}{\texttt{NEGF}}\par
\noindent\texttt{\$D8}\quad \hyperlink{l:NEGI}{\texttt{NEGI}}\par
\noindent\texttt{\$D9}\quad \hyperlink{l:FTOI}{\texttt{FTOI}}\par
\noindent\texttt{\$DA}\quad \hyperlink{l:ITOF}{\texttt{ITOF}}\par
\noindent\texttt{\$DB}\quad \hyperlink{l:SKPARG}{\texttt{SKPARG}}\par
\noindent\texttt{\$DC}\quad \hyperlink{l:OPEQV}{\texttt{OPEQV}}\hfill\texttt{EQV}\par
\noindent\texttt{\$DD}\quad \hyperlink{l:OPIMP}{\texttt{OPIMP}}\hfill\texttt{IMP}\par
\noindent\texttt{\$DE}\quad \hyperlink{l:OPOR}{\texttt{OPOR}}\hfill\texttt{OR}\par
\noindent\texttt{\$DF}\quad \hyperlink{l:OPXOR}{\texttt{OPXOR}}\hfill\texttt{XOR}\par
\noindent\texttt{\$E0}\quad \hyperlink{l:OPAND}{\texttt{OPAND}}\hfill\texttt{AND}\par
\noindent\texttt{\$E1}\quad \hyperlink{l:OPNOT}{\texttt{OPNOT}}\hfill\texttt{NOT}\par
\noindent\texttt{\$E2}\quad \hyperlink{l:EQF}{\texttt{EQF}}\hfill\texttt{=}\par
\noindent\texttt{\$E3}\quad \hyperlink{l:EQI}{\texttt{EQI}}\hfill\texttt{=}\par
\noindent\texttt{\$E4}\quad \hyperlink{l:EQS}{\texttt{EQS}}\hfill\texttt{=}\par
\noindent\texttt{\$E5}\quad \hyperlink{l:NEF}{\texttt{NEF}}\hfill\texttt{{\char60}{\char62}}\par
\noindent\texttt{\$E6}\quad \hyperlink{l:NEI}{\texttt{NEI}}\hfill\texttt{{\char60}{\char62}}\par
\noindent\texttt{\$E7}\quad \hyperlink{l:NES}{\texttt{NES}}\hfill\texttt{{\char60}{\char62}}\par
\noindent\texttt{\$E8}\quad \hyperlink{l:LTF}{\texttt{LTF}}\hfill\texttt{{\char60}}\par
\noindent\texttt{\$E9}\quad \hyperlink{l:LTI}{\texttt{LTI}}\hfill\texttt{{\char60}}\par
\noindent\texttt{\$EA}\quad \hyperlink{l:LTS}{\texttt{LTS}}\hfill\texttt{{\char60}}\par
\noindent\texttt{\$EB}\quad \hyperlink{l:GEF}{\texttt{GEF}}\hfill\texttt{{\char62}=}\par
\noindent\texttt{\$EC}\quad \hyperlink{l:GEI}{\texttt{GEI}}\hfill\texttt{{\char62}=}\par
\noindent\texttt{\$ED}\quad \hyperlink{l:GES}{\texttt{GES}}\hfill\texttt{{\char62}=}\par
\noindent\texttt{\$EE}\quad \hyperlink{l:GTF}{\texttt{GTF}}\hfill\texttt{{\char62}}\par
\noindent\texttt{\$EF}\quad \hyperlink{l:GTI}{\texttt{GTI}}\hfill\texttt{{\char62}}\par
\noindent\texttt{\$F0}\quad \hyperlink{l:GTS}{\texttt{GTS}}\hfill\texttt{{\char62}}\par
\noindent\texttt{\$F1}\quad \hyperlink{l:LEF}{\texttt{LEF}}\hfill\texttt{{\char60}=}\par
\noindent\texttt{\$F2}\quad \hyperlink{l:LEI}{\texttt{LEI}}\hfill\texttt{{\char60}=}\par
\noindent\texttt{\$F3}\quad \hyperlink{l:LES}{\texttt{LES}}\hfill\texttt{{\char60}=}\par
\noindent\texttt{\$F4}\quad \hyperlink{l:ADDF}{\texttt{ADDF}}\hfill\texttt{+}\par
\noindent\texttt{\$F5}\quad \hyperlink{l:ADDI}{\texttt{ADDI}}\hfill\texttt{+}\par
\noindent\texttt{\$F6}\quad \hyperlink{l:CONCAT}{\texttt{CONCAT}}\hfill\texttt{+}\par
\noindent\texttt{\$F7}\quad \hyperlink{l:SUBF}{\texttt{SUBF}}\hfill\texttt{-}\par
\noindent\texttt{\$F8}\quad \hyperlink{l:SUBI}{\texttt{SUBI}}\hfill\texttt{-}\par
\noindent\texttt{\$F9}\quad \hyperlink{l:MULF}{\texttt{MULF}}\hfill\texttt{*}\par
\noindent\texttt{\$FA}\quad \hyperlink{l:MULI}{\texttt{MULI}}\hfill\texttt{*}\par
\noindent\texttt{\$FB}\quad \hyperlink{l:DIVF}{\texttt{DIVF}}\hfill\texttt{/}\par
\noindent\texttt{\$FC}\quad \hyperlink{l:DIVI}{\texttt{DIVI}}\hfill\texttt{/}\par
\noindent\texttt{\$FD}\quad \hyperlink{l:POWF}{\texttt{POWF}}\hfill\texttt{\^{}}\par
\noindent\texttt{\$FE}\quad \hyperlink{l:POWI}{\texttt{POWI}}\hfill\texttt{\^{}}\par
\noindent\texttt{\$FF}\quad \hyperlink{l:POWFI}{\texttt{POWFI}}\hfill\texttt{\^{}}\par
\end{multicols}

% kapitel/symboler.tex -- fryst ur tex80.py 2026-10-06 (omstruktureringen steg 5);
% rättas här för hand.

\chapter{Symbolförteckning}
\noindent Förteckningen är till för att hitta ett namn som står i texten
när man inte vet adressen. Kapitlet Rutinerna (i boken)
har namnen i adressordning med uppgiften i korthet; här står samma namn i
bokstavsordning. Varje rad ger namnet, adressen och sidan i Rutinerna där
namnet står. I den elektroniska utgåvan är namnet dessutom en länk dit.
Namnen är valda för den här boken; ROM:en själv innehåller inga.
\par\smallskip\noindent Namn som bara är en bokstav och en adress tas
inte med, eftersom namnet redan säger var det står i Rutinerna. Ett J är
en ingång i början av ROM:en och ett D en text eller en tabell, till
exempel \namn{D0501} på \adr{0501}. Hoppmålen inne i rutinerna, ett L
och adressen, står inte heller i Rutinerna. Variablerna i RAM och
adresserna utanför ROM:en står i kapitlet Systemvariabler och externa
adresser (i boken).
\par\medskip
\begin{multicols}{3}\footnotesize\raggedright
\noindent\hyperlink{l:AADD}{\texttt{AADD}}\hfill\texttt{\$37A6~\phantom{000}\llap{\pageref{l:AADD}}}\par
\noindent\hyperlink{l:ACOMP}{\texttt{ACOMP}}\hfill\texttt{\$37F6~\phantom{000}\llap{\pageref{l:ACOMP}}}\par
\noindent\hyperlink{l:ADDF}{\texttt{ADDF}}\hfill\texttt{\$2AB7~\phantom{000}\llap{\pageref{l:ADDF}}}\par
\noindent\hyperlink{l:ADDI}{\texttt{ADDI}}\hfill\texttt{\$2B5E~\phantom{000}\llap{\pageref{l:ADDI}}}\par
\noindent\hyperlink{l:ADIV}{\texttt{ADIV}}\hfill\texttt{\$383B~\phantom{000}\llap{\pageref{l:ADIV}}}\par
\noindent\hyperlink{l:ALLOC}{\texttt{ALLOC}}\hfill\texttt{\$3374~\phantom{000}\llap{\pageref{l:ALLOC}}}\par
\noindent\hyperlink{l:AMUL}{\texttt{AMUL}}\hfill\texttt{\$38B8~\phantom{000}\llap{\pageref{l:AMUL}}}\par
\noindent\hyperlink{l:ARGWLK}{\texttt{ARGWLK}}\hfill\texttt{\$33B0~\phantom{000}\llap{\pageref{l:ARGWLK}}}\par
\noindent\hyperlink{l:ASADDB}{\texttt{ASADDB}}\hfill\texttt{\$3ABA~\phantom{000}\llap{\pageref{l:ASADDB}}}\par
\noindent\hyperlink{l:ASADDN}{\texttt{ASADDN}}\hfill\texttt{\$390D~\phantom{000}\llap{\pageref{l:ASADDN}}}\par
\noindent\hyperlink{l:ASALG1}{\texttt{ASALG1}}\hfill\texttt{\$3A32~\phantom{000}\llap{\pageref{l:ASALG1}}}\par
\noindent\hyperlink{l:ASALGN}{\texttt{ASALGN}}\hfill\texttt{\$3A12~\phantom{000}\llap{\pageref{l:ASALGN}}}\par
\noindent\hyperlink{l:ASALST}{\texttt{ASALST}}\hfill\texttt{\$3B72~\phantom{000}\llap{\pageref{l:ASALST}}}\par
\noindent\hyperlink{l:ASBAD}{\texttt{ASBAD}}\hfill\texttt{\$39BF~\phantom{000}\llap{\pageref{l:ASBAD}}}\par
\noindent\hyperlink{l:ASBLST}{\texttt{ASBLST}}\hfill\texttt{\$3B6D~\phantom{000}\llap{\pageref{l:ASBLST}}}\par
\noindent\hyperlink{l:ASBORN}{\texttt{ASBORN}}\hfill\texttt{\$3948~\phantom{000}\llap{\pageref{l:ASBORN}}}\par
\noindent\hyperlink{l:ASCAR}{\texttt{ASCAR}}\hfill\texttt{\$2BBA~\phantom{000}\llap{\pageref{l:ASCAR}}}\par
\noindent\hyperlink{l:ASCARG}{\texttt{ASCARG}}\hfill\texttt{\$2C6C~\phantom{000}\llap{\pageref{l:ASCARG}}}\par
\noindent\hyperlink{l:ASCARN}{\texttt{ASCARN}}\hfill\texttt{\$391A~\phantom{000}\llap{\pageref{l:ASCARN}}}\par
\noindent\hyperlink{l:ASCCMP}{\texttt{ASCCMP}}\hfill\texttt{\$2C51~\phantom{000}\llap{\pageref{l:ASCCMP}}}\par
\noindent\hyperlink{l:ASCERR}{\texttt{ASCERR}}\hfill\texttt{\$2C4B~\phantom{000}\llap{\pageref{l:ASCERR}}}\par
\noindent\hyperlink{l:ASCMPS}{\texttt{ASCMPS}}\hfill\texttt{\$3817~\phantom{000}\llap{\pageref{l:ASCMPS}}}\par
\noindent\hyperlink{l:ASCNT}{\texttt{ASCNT}}\hfill\texttt{\$3B68~\phantom{000}\llap{\pageref{l:ASCNT}}}\par
\noindent\hyperlink{l:ASDIG}{\texttt{ASDIG}}\hfill\texttt{\$3B4D~\phantom{000}\llap{\pageref{l:ASDIG}}}\par
\noindent\hyperlink{l:ASDIG2}{\texttt{ASDIG2}}\hfill\texttt{\$395E~\phantom{000}\llap{\pageref{l:ASDIG2}}}\par
\noindent\hyperlink{l:ASDLP}{\texttt{ASDLP}}\hfill\texttt{\$3868~\phantom{000}\llap{\pageref{l:ASDLP}}}\par
\noindent\hyperlink{l:ASDOT}{\texttt{ASDOT}}\hfill\texttt{\$3B47~\phantom{000}\llap{\pageref{l:ASDOT}}}\par
\noindent\hyperlink{l:ASENTR}{\texttt{ASENTR}}\hfill\texttt{\$396F~\phantom{000}\llap{\pageref{l:ASENTR}}}\par
\noindent\hyperlink{l:ASFIN}{\texttt{ASFIN}}\hfill\texttt{\$3B00~\phantom{000}\llap{\pageref{l:ASFIN}}}\par
\noindent\hyperlink{l:ASGNF}{\texttt{ASGNF}}\hfill\texttt{\$2989~\phantom{000}\llap{\pageref{l:ASGNF}}}\par
\noindent\hyperlink{l:ASGNI}{\texttt{ASGNI}}\hfill\texttt{\$29B4~\phantom{000}\llap{\pageref{l:ASGNI}}}\par
\noindent\hyperlink{l:ASGNS}{\texttt{ASGNS}}\hfill\texttt{\$29C3~\phantom{000}\llap{\pageref{l:ASGNS}}}\par
\noindent\hyperlink{l:ASLEN}{\texttt{ASLEN}}\hfill\texttt{\$38FC~\phantom{000}\llap{\pageref{l:ASLEN}}}\par
\noindent\hyperlink{l:ASLEN1}{\texttt{ASLEN1}}\hfill\texttt{\$3902~\phantom{000}\llap{\pageref{l:ASLEN1}}}\par
\noindent\hyperlink{l:ASMADD}{\texttt{ASMADD}}\hfill\texttt{\$37B4~\phantom{000}\llap{\pageref{l:ASMADD}}}\par
\noindent\hyperlink{l:ASMLP}{\texttt{ASMLP}}\hfill\texttt{\$38CC~\phantom{000}\llap{\pageref{l:ASMLP}}}\par
\noindent\hyperlink{l:ASMSUB}{\texttt{ASMSUB}}\hfill\texttt{\$37C9~\phantom{000}\llap{\pageref{l:ASMSUB}}}\par
\noindent\hyperlink{l:ASNEG}{\texttt{ASNEG}}\hfill\texttt{\$3AD6~\phantom{000}\llap{\pageref{l:ASNEG}}}\par
\noindent\hyperlink{l:ASNZ}{\texttt{ASNZ}}\hfill\texttt{\$3B53~\phantom{000}\llap{\pageref{l:ASNZ}}}\par
\noindent\hyperlink{l:ASPARS}{\texttt{ASPARS}}\hfill\texttt{\$39C5~\phantom{000}\llap{\pageref{l:ASPARS}}}\par
\noindent\hyperlink{l:ASPUT}{\texttt{ASPUT}}\hfill\texttt{\$3926~\phantom{000}\llap{\pageref{l:ASPUT}}}\par
\noindent\hyperlink{l:ASPUTB}{\texttt{ASPUTB}}\hfill\texttt{\$3954~\phantom{000}\llap{\pageref{l:ASPUTB}}}\par
\noindent\hyperlink{l:ASREST}{\texttt{ASREST}}\hfill\texttt{\$3ACD~\phantom{000}\llap{\pageref{l:ASREST}}}\par
\noindent\hyperlink{l:ASRET}{\texttt{ASRET}}\hfill\texttt{\$3B3E~\phantom{000}\llap{\pageref{l:ASRET}}}\par
\noindent\hyperlink{l:ASSETB}{\texttt{ASSETB}}\hfill\texttt{\$38AB~\phantom{000}\llap{\pageref{l:ASSETB}}}\par
\noindent\hyperlink{l:ASSUBN}{\texttt{ASSUBN}}\hfill\texttt{\$393B~\phantom{000}\llap{\pageref{l:ASSUBN}}}\par
\noindent\hyperlink{l:ASUB}{\texttt{ASUB}}\hfill\texttt{\$37AE~\phantom{000}\llap{\pageref{l:ASUB}}}\par
\noindent\hyperlink{l:ASWIDE}{\texttt{ASWIDE}}\hfill\texttt{\$3A67~\phantom{000}\llap{\pageref{l:ASWIDE}}}\par
\noindent\hyperlink{l:ASWORK}{\texttt{ASWORK}}\hfill\texttt{\$3A9D~\phantom{000}\llap{\pageref{l:ASWORK}}}\par
\noindent\hyperlink{l:ASZERA}{\texttt{ASZERA}}\hfill\texttt{\$37E6~\phantom{000}\llap{\pageref{l:ASZERA}}}\par
\noindent\hyperlink{l:ATNPOS}{\texttt{ATNPOS}}\hfill\texttt{\$3B8A~\phantom{000}\llap{\pageref{l:ATNPOS}}}\par
\noindent\hyperlink{l:ATNSER}{\texttt{ATNSER}}\hfill\texttt{\$3C02~\phantom{000}\llap{\pageref{l:ATNSER}}}\par
\noindent\hyperlink{l:BRKSTP}{\texttt{BRKSTP}}\hfill\texttt{\$2323~\phantom{000}\llap{\pageref{l:BRKSTP}}}\par
\noindent\hyperlink{l:CALLBC}{\texttt{CALLBC}}\hfill\texttt{\$2C47~\phantom{000}\llap{\pageref{l:CALLBC}}}\par
\noindent\hyperlink{l:CASBRS}{\texttt{CASBRS}}\hfill\texttt{\$03F4~\phantom{000}\llap{\pageref{l:CASBRS}}}\par
\noindent\hyperlink{l:CASCLO}{\texttt{CASCLO}}\hfill\texttt{\$0428~\phantom{000}\llap{\pageref{l:CASCLO}}}\par
\noindent\hyperlink{l:CASGET}{\texttt{CASGET}}\hfill\texttt{\$0509~\phantom{000}\llap{\pageref{l:CASGET}}}\par
\noindent\hyperlink{l:CASHLF}{\texttt{CASHLF}}\hfill\texttt{\$0484~\phantom{000}\llap{\pageref{l:CASHLF}}}\par
\noindent\hyperlink{l:CASINT}{\texttt{CASINT}}\hfill\texttt{\$0594~\phantom{000}\llap{\pageref{l:CASINT}}}\par
\noindent\hyperlink{l:CASOPN}{\texttt{CASOPN}}\hfill\texttt{\$0491~\phantom{000}\llap{\pageref{l:CASOPN}}}\par
\noindent\hyperlink{l:CASPRE}{\texttt{CASPRE}}\hfill\texttt{\$03BF~\phantom{000}\llap{\pageref{l:CASPRE}}}\par
\noindent\hyperlink{l:CASPUT}{\texttt{CASPUT}}\hfill\texttt{\$040E~\phantom{000}\llap{\pageref{l:CASPUT}}}\par
\noindent\hyperlink{l:CASRBL}{\texttt{CASRBL}}\hfill\texttt{\$052B~\phantom{000}\llap{\pageref{l:CASRBL}}}\par
\noindent\hyperlink{l:CASTAB}{\texttt{CASTAB}}\hfill\texttt{\$03A4~\phantom{000}\llap{\pageref{l:CASTAB}}}\par
\noindent\hyperlink{l:CASW0}{\texttt{CASW0}}\hfill\texttt{\$0476~\phantom{000}\llap{\pageref{l:CASW0}}}\par
\noindent\hyperlink{l:CASWBI}{\texttt{CASWBI}}\hfill\texttt{\$0477~\phantom{000}\llap{\pageref{l:CASWBI}}}\par
\noindent\hyperlink{l:CASWBL}{\texttt{CASWBL}}\hfill\texttt{\$0439~\phantom{000}\llap{\pageref{l:CASWBL}}}\par
\noindent\hyperlink{l:CASWBY}{\texttt{CASWBY}}\hfill\texttt{\$046A~\phantom{000}\llap{\pageref{l:CASWBY}}}\par
\noindent\hyperlink{l:CHAIN}{\texttt{CHAIN}}\hfill\texttt{\$0D3F~\phantom{000}\llap{\pageref{l:CHAIN}}}\par
\noindent\hyperlink{l:CLEAR}{\texttt{CLEAR}}\hfill\texttt{\$0A80~\phantom{000}\llap{\pageref{l:CLEAR}}}\par
\noindent\hyperlink{l:CLOSE}{\texttt{CLOSE}}\hfill\texttt{\$087F~\phantom{000}\llap{\pageref{l:CLOSE}}}\par
\noindent\hyperlink{l:CLOSST}{\texttt{CLOSST}}\hfill\texttt{\$20BF~\phantom{000}\llap{\pageref{l:CLOSST}}}\par
\noindent\hyperlink{l:CLRDOT}{\texttt{CLRDOT}}\hfill\texttt{\$22CF~\phantom{000}\llap{\pageref{l:CLRDOT}}}\par
\noindent\hyperlink{l:CLRROW}{\texttt{CLRROW}}\hfill\texttt{\$0267~\phantom{000}\llap{\pageref{l:CLRROW}}}\par
\noindent\hyperlink{l:CLS}{\texttt{CLS}}\hfill\texttt{\$0276~\phantom{000}\llap{\pageref{l:CLS}}}\par
\noindent\hyperlink{l:CLSALL}{\texttt{CLSALL}}\hfill\texttt{\$20CD~\phantom{000}\llap{\pageref{l:CLSALL}}}\par
\noindent\hyperlink{l:CMDED}{\texttt{CMDED}}\hfill\texttt{\$0AA1~\phantom{000}\llap{\pageref{l:CMDED}}}\par
\noindent\hyperlink{l:CMDKW}{\texttt{CMDKW}}\hfill\texttt{\$0A20~\phantom{000}\llap{\pageref{l:CMDKW}}}\par
\noindent\hyperlink{l:CMDLIST}{\texttt{CMDLIST}}\hfill\texttt{\$0AEF~\phantom{000}\llap{\pageref{l:CMDLIST}}}\par
\noindent\hyperlink{l:CMDLOAD}{\texttt{CMDLOAD}}\hfill\texttt{\$0BA8~\phantom{000}\llap{\pageref{l:CMDLOAD}}}\par
\noindent\hyperlink{l:CMDLP}{\texttt{CMDLP}}\hfill\texttt{\$00DB~\phantom{000}\llap{\pageref{l:CMDLP}}}\par
\noindent\hyperlink{l:CMDMERG}{\texttt{CMDMERG}}\hfill\texttt{\$0BA9~\phantom{000}\llap{\pageref{l:CMDMERG}}}\par
\noindent\hyperlink{l:CMDREN}{\texttt{CMDREN}}\hfill\texttt{\$0C8A~\phantom{000}\llap{\pageref{l:CMDREN}}}\par
\noindent\hyperlink{l:CMDRUN}{\texttt{CMDRUN}}\hfill\texttt{\$0D4C~\phantom{000}\llap{\pageref{l:CMDRUN}}}\par
\noindent\hyperlink{l:CMDSAVE}{\texttt{CMDSAVE}}\hfill\texttt{\$0DC6~\phantom{000}\llap{\pageref{l:CMDSAVE}}}\par
\noindent\hyperlink{l:CMDTAB}{\texttt{CMDTAB}}\hfill\texttt{\$0A5A~\phantom{000}\llap{\pageref{l:CMDTAB}}}\par
\noindent\hyperlink{l:CMPADD}{\texttt{CMPADD}}\hfill\texttt{\$27FB~\phantom{000}\llap{\pageref{l:CMPADD}}}\par
\noindent\hyperlink{l:CMPARG}{\texttt{CMPARG}}\hfill\texttt{\$264E~\phantom{000}\llap{\pageref{l:CMPARG}}}\par
\noindent\hyperlink{l:CMPB9}{\texttt{CMPB9}}\hfill\texttt{\$1CAF~\phantom{000}\llap{\pageref{l:CMPB9}}}\par
\noindent\hyperlink{l:CMPBF}{\texttt{CMPBF}}\hfill\texttt{\$1BB8~\phantom{000}\llap{\pageref{l:CMPBF}}}\par
\noindent\hyperlink{l:CMPCLS}{\texttt{CMPCLS}}\hfill\texttt{\$2071~\phantom{000}\llap{\pageref{l:CMPCLS}}}\par
\noindent\hyperlink{l:CMPDEF}{\texttt{CMPDEF}}\hfill\texttt{\$1A62~\phantom{000}\llap{\pageref{l:CMPDEF}}}\par
\noindent\hyperlink{l:CMPDIM}{\texttt{CMPDIM}}\hfill\texttt{\$1BAA~\phantom{000}\llap{\pageref{l:CMPDIM}}}\par
\noindent\hyperlink{l:CMPDMI}{\texttt{CMPDMI}}\hfill\texttt{\$1BBD~\phantom{000}\llap{\pageref{l:CMPDMI}}}\par
\noindent\hyperlink{l:CMPDOT}{\texttt{CMPDOT}}\hfill\texttt{\$22B6~\phantom{000}\llap{\pageref{l:CMPDOT}}}\par
\noindent\hyperlink{l:CMPERR}{\texttt{CMPERR}}\hfill\texttt{\$1BF7~\phantom{000}\llap{\pageref{l:CMPERR}}}\par
\noindent\hyperlink{l:CMPEXB}{\texttt{CMPEXB}}\hfill\texttt{\$2362~\phantom{000}\llap{\pageref{l:CMPEXB}}}\par
\noindent\hyperlink{l:CMPEXP}{\texttt{CMPEXP}}\hfill\texttt{\$2701~\phantom{000}\llap{\pageref{l:CMPEXP}}}\par
\noindent\hyperlink{l:CMPFLV}{\texttt{CMPFLV}}\hfill\texttt{\$1BA5~\phantom{000}\llap{\pageref{l:CMPFLV}}}\par
\noindent\hyperlink{l:CMPFN}{\texttt{CMPFN}}\hfill\texttt{\$1AB8~\phantom{000}\llap{\pageref{l:CMPFN}}}\par
\noindent\hyperlink{l:CMPFNC}{\texttt{CMPFNC}}\hfill\texttt{\$261C~\phantom{000}\llap{\pageref{l:CMPFNC}}}\par
\noindent\hyperlink{l:CMPFOR}{\texttt{CMPFOR}}\hfill\texttt{\$1C83~\phantom{000}\llap{\pageref{l:CMPFOR}}}\par
\noindent\hyperlink{l:CMPGET}{\texttt{CMPGET}}\hfill\texttt{\$23A8~\phantom{000}\llap{\pageref{l:CMPGET}}}\par
\noindent\hyperlink{l:CMPIF}{\texttt{CMPIF}}\hfill\texttt{\$1DFB~\phantom{000}\llap{\pageref{l:CMPIF}}}\par
\noindent\hyperlink{l:CMPINL}{\texttt{CMPINL}}\hfill\texttt{\$1F6B~\phantom{000}\llap{\pageref{l:CMPINL}}}\par
\noindent\hyperlink{l:CMPINP}{\texttt{CMPINP}}\hfill\texttt{\$1E58~\phantom{000}\llap{\pageref{l:CMPINP}}}\par
\noindent\hyperlink{l:CMPINX}{\texttt{CMPINX}}\hfill\texttt{\$2370~\phantom{000}\llap{\pageref{l:CMPINX}}}\par
\noindent\hyperlink{l:CMPIVR}{\texttt{CMPIVR}}\hfill\texttt{\$239F~\phantom{000}\llap{\pageref{l:CMPIVR}}}\par
\noindent\hyperlink{l:CMPLET}{\texttt{CMPLET}}\hfill\texttt{\$1B89~\phantom{000}\llap{\pageref{l:CMPLET}}}\par
\noindent\hyperlink{l:CMPNOP}{\texttt{CMPNOP}}\hfill\texttt{\$226A~\phantom{000}\llap{\pageref{l:CMPNOP}}}\par
\noindent\hyperlink{l:CMPNOT}{\texttt{CMPNOT}}\hfill\texttt{\$2826~\phantom{000}\llap{\pageref{l:CMPNOT}}}\par
\noindent\hyperlink{l:CMPNUM}{\texttt{CMPNUM}}\hfill\texttt{\$2689~\phantom{000}\llap{\pageref{l:CMPNUM}}}\par
\noindent\hyperlink{l:CMPNXT}{\texttt{CMPNXT}}\hfill\texttt{\$1CB7~\phantom{000}\llap{\pageref{l:CMPNXT}}}\par
\noindent\hyperlink{l:CMPON}{\texttt{CMPON}}\hfill\texttt{\$1FBF~\phantom{000}\llap{\pageref{l:CMPON}}}\par
\noindent\hyperlink{l:CMPOPN}{\texttt{CMPOPN}}\hfill\texttt{\$205D~\phantom{000}\llap{\pageref{l:CMPOPN}}}\par
\noindent\hyperlink{l:CMPOUT}{\texttt{CMPOUT}}\hfill\texttt{\$20FF~\phantom{000}\llap{\pageref{l:CMPOUT}}}\par
\noindent\hyperlink{l:CMPPOK}{\texttt{CMPPOK}}\hfill\texttt{\$2104~\phantom{000}\llap{\pageref{l:CMPPOK}}}\par
\noindent\hyperlink{l:CMPPRT}{\texttt{CMPPRT}}\hfill\texttt{\$2154~\phantom{000}\llap{\pageref{l:CMPPRT}}}\par
\noindent\hyperlink{l:CMPRD}{\texttt{CMPRD}}\hfill\texttt{\$1E6E~\phantom{000}\llap{\pageref{l:CMPRD}}}\par
\noindent\hyperlink{l:CMPREF}{\texttt{CMPREF}}\hfill\texttt{\$1DC7~\phantom{000}\llap{\pageref{l:CMPREF}}}\par
\noindent\hyperlink{l:CMPREM}{\texttt{CMPREM}}\hfill\texttt{\$226B~\phantom{000}\llap{\pageref{l:CMPREM}}}\par
\noindent\hyperlink{l:CMPRES}{\texttt{CMPRES}}\hfill\texttt{\$2281~\phantom{000}\llap{\pageref{l:CMPRES}}}\par
\noindent\hyperlink{l:CMPSTR}{\texttt{CMPSTR}}\hfill\texttt{\$25C6~\phantom{000}\llap{\pageref{l:CMPSTR}}}\par
\noindent\hyperlink{l:CMPSTX}{\texttt{CMPSTX}}\hfill\texttt{\$236C~\phantom{000}\llap{\pageref{l:CMPSTX}}}\par
\noindent\hyperlink{l:CMPSV}{\texttt{CMPSV}}\hfill\texttt{\$237F~\phantom{000}\llap{\pageref{l:CMPSV}}}\par
\noindent\hyperlink{l:CMPSVR}{\texttt{CMPSVR}}\hfill\texttt{\$2392~\phantom{000}\llap{\pageref{l:CMPSVR}}}\par
\noindent\hyperlink{l:CMPTYX}{\texttt{CMPTYX}}\hfill\texttt{\$2372~\phantom{000}\llap{\pageref{l:CMPTYX}}}\par
\noindent\hyperlink{l:CMPVAR}{\texttt{CMPVAR}}\hfill\texttt{\$25EE~\phantom{000}\llap{\pageref{l:CMPVAR}}}\par
\noindent\hyperlink{l:CMPVR}{\texttt{CMPVR}}\hfill\texttt{\$2397~\phantom{000}\llap{\pageref{l:CMPVR}}}\par
\noindent\hyperlink{l:COLD}{\texttt{COLD}}\hfill\texttt{\$00C6~\phantom{000}\llap{\pageref{l:COLD}}}\par
\noindent\hyperlink{l:CONCAT}{\texttt{CONCAT}}\hfill\texttt{\$2F47~\phantom{000}\llap{\pageref{l:CONCAT}}}\par
\noindent\hyperlink{l:CONSTF}{\texttt{CONSTF}}\hfill\texttt{\$2A36~\phantom{000}\llap{\pageref{l:CONSTF}}}\par
\noindent\hyperlink{l:CONSTI}{\texttt{CONSTI}}\hfill\texttt{\$2AE3~\phantom{000}\llap{\pageref{l:CONSTI}}}\par
\noindent\hyperlink{l:CONSTS}{\texttt{CONSTS}}\hfill\texttt{\$2ED7~\phantom{000}\llap{\pageref{l:CONSTS}}}\par
\noindent\hyperlink{l:COSSER}{\texttt{COSSER}}\hfill\texttt{\$3FF4~\phantom{000}\llap{\pageref{l:COSSER}}}\par
\noindent\hyperlink{l:CRLF}{\texttt{CRLF}}\hfill\texttt{\$0234~\phantom{000}\llap{\pageref{l:CRLF}}}\par
\noindent\hyperlink{l:CURADR}{\texttt{CURADR}}\hfill\texttt{\$0293~\phantom{000}\llap{\pageref{l:CURADR}}}\par
\noindent\hyperlink{l:CURFIL}{\texttt{CURFIL}}\hfill\texttt{\$08FC~\phantom{000}\llap{\pageref{l:CURFIL}}}\par
\noindent\hyperlink{l:DATA}{\texttt{DATA}}\hfill\texttt{\$2273~\phantom{000}\llap{\pageref{l:DATA}}}\par
\noindent\hyperlink{l:DEF}{\texttt{DEF}}\hfill\texttt{\$1AAA~\phantom{000}\llap{\pageref{l:DEF}}}\par
\noindent\hyperlink{l:DETEXP}{\texttt{DETEXP}}\hfill\texttt{\$1258~\phantom{000}\llap{\pageref{l:DETEXP}}}\par
\noindent\hyperlink{l:DETOK}{\texttt{DETOK}}\hfill\texttt{\$1195~\phantom{000}\llap{\pageref{l:DETOK}}}\par
\noindent\hyperlink{l:DETOP}{\texttt{DETOP}}\hfill\texttt{\$1398~\phantom{000}\llap{\pageref{l:DETOP}}}\par
\noindent\hyperlink{l:DETTAB}{\texttt{DETTAB}}\hfill\texttt{\$12B6~\phantom{000}\llap{\pageref{l:DETTAB}}}\par
\noindent\hyperlink{l:DETVAR}{\texttt{DETVAR}}\hfill\texttt{\$1422~\phantom{000}\llap{\pageref{l:DETVAR}}}\par
\noindent\hyperlink{l:DEVDFL}{\texttt{DEVDFL}}\hfill\texttt{\$0193~\phantom{000}\llap{\pageref{l:DEVDFL}}}\par
\noindent\hyperlink{l:DEVIEC}{\texttt{DEVIEC}}\hfill\texttt{\$01A1~\phantom{000}\llap{\pageref{l:DEVIEC}}}\par
\noindent\hyperlink{l:DEVLST}{\texttt{DEVLST}}\hfill\texttt{\$018C~\phantom{000}\llap{\pageref{l:DEVLST}}}\par
\noindent\hyperlink{l:DEVPR}{\texttt{DEVPR}}\hfill\texttt{\$019A~\phantom{000}\llap{\pageref{l:DEVPR}}}\par
\noindent\hyperlink{l:DIM}{\texttt{DIM}}\hfill\texttt{\$1BFA~\phantom{000}\llap{\pageref{l:DIM}}}\par
\noindent\hyperlink{l:DIM1}{\texttt{DIM1}}\hfill\texttt{\$32E6~\phantom{000}\llap{\pageref{l:DIM1}}}\par
\noindent\hyperlink{l:DIM2}{\texttt{DIM2}}\hfill\texttt{\$3311~\phantom{000}\llap{\pageref{l:DIM2}}}\par
\noindent\hyperlink{l:DIMELM}{\texttt{DIMELM}}\hfill\texttt{\$1BFE~\phantom{000}\llap{\pageref{l:DIMELM}}}\par
\noindent\hyperlink{l:DIMSTR}{\texttt{DIMSTR}}\hfill\texttt{\$3338~\phantom{000}\llap{\pageref{l:DIMSTR}}}\par
\noindent\hyperlink{l:DIVF}{\texttt{DIVF}}\hfill\texttt{\$2AC9~\phantom{000}\llap{\pageref{l:DIVF}}}\par
\noindent\hyperlink{l:DIVI}{\texttt{DIVI}}\hfill\texttt{\$2B73~\phantom{000}\llap{\pageref{l:DIVI}}}\par
\noindent\hyperlink{l:DOCMD}{\texttt{DOCMD}}\hfill\texttt{\$0133~\phantom{000}\llap{\pageref{l:DOCMD}}}\par
\noindent\hyperlink{l:DOTADR}{\texttt{DOTADR}}\hfill\texttt{\$22EE~\phantom{000}\llap{\pageref{l:DOTADR}}}\par
\noindent\hyperlink{l:DOTARG}{\texttt{DOTARG}}\hfill\texttt{\$22E3~\phantom{000}\llap{\pageref{l:DOTARG}}}\par
\noindent\hyperlink{l:DVCALL}{\texttt{DVCALL}}\hfill\texttt{\$08A5~\phantom{000}\llap{\pageref{l:DVCALL}}}\par
\noindent\hyperlink{l:DVCLOS}{\texttt{DVCLOS}}\hfill\texttt{\$0890~\phantom{000}\llap{\pageref{l:DVCLOS}}}\par
\noindent\hyperlink{l:DVGET}{\texttt{DVGET}}\hfill\texttt{\$0899~\phantom{000}\llap{\pageref{l:DVGET}}}\par
\noindent\hyperlink{l:DVINP}{\texttt{DVINP}}\hfill\texttt{\$0893~\phantom{000}\llap{\pageref{l:DVINP}}}\par
\noindent\hyperlink{l:DVOPEN}{\texttt{DVOPEN}}\hfill\texttt{\$088A~\phantom{000}\llap{\pageref{l:DVOPEN}}}\par
\noindent\hyperlink{l:DVPREP}{\texttt{DVPREP}}\hfill\texttt{\$088D~\phantom{000}\llap{\pageref{l:DVPREP}}}\par
\noindent\hyperlink{l:DVPRT}{\texttt{DVPRT}}\hfill\texttt{\$0896~\phantom{000}\llap{\pageref{l:DVPRT}}}\par
\noindent\hyperlink{l:DVPUT}{\texttt{DVPUT}}\hfill\texttt{\$089C~\phantom{000}\llap{\pageref{l:DVPUT}}}\par
\noindent\hyperlink{l:EDTAB}{\texttt{EDTAB}}\hfill\texttt{\$0348~\phantom{000}\llap{\pageref{l:EDTAB}}}\par
\noindent\hyperlink{l:ELMADR}{\texttt{ELMADR}}\hfill\texttt{\$2CFA~\phantom{000}\llap{\pageref{l:ELMADR}}}\par
\noindent\hyperlink{l:ELMSIZ}{\texttt{ELMSIZ}}\hfill\texttt{\$1C6C~\phantom{000}\llap{\pageref{l:ELMSIZ}}}\par
\noindent\hyperlink{l:ENDDIM}{\texttt{ENDDIM}}\hfill\texttt{\$1C7E~\phantom{000}\llap{\pageref{l:ENDDIM}}}\par
\noindent\hyperlink{l:ENDF}{\texttt{ENDF}}\hfill\texttt{\$285F~\phantom{000}\llap{\pageref{l:ENDF}}}\par
\noindent\hyperlink{l:ENDPRG}{\texttt{ENDPRG}}\hfill\texttt{\$00C9~\phantom{000}\llap{\pageref{l:ENDPRG}}}\par
\noindent\hyperlink{l:ENDS}{\texttt{ENDS}}\hfill\texttt{\$2877~\phantom{000}\llap{\pageref{l:ENDS}}}\par
\noindent\hyperlink{l:ENDVF}{\texttt{ENDVF}}\hfill\texttt{\$2884~\phantom{000}\llap{\pageref{l:ENDVF}}}\par
\noindent\hyperlink{l:ENDVS}{\texttt{ENDVS}}\hfill\texttt{\$2887~\phantom{000}\llap{\pageref{l:ENDVS}}}\par
\noindent\hyperlink{l:EQF}{\texttt{EQF}}\hfill\texttt{\$2A4C~\phantom{000}\llap{\pageref{l:EQF}}}\par
\noindent\hyperlink{l:EQI}{\texttt{EQI}}\hfill\texttt{\$2B2A~\phantom{000}\llap{\pageref{l:EQI}}}\par
\noindent\hyperlink{l:EQS}{\texttt{EQS}}\hfill\texttt{\$2FE4~\phantom{000}\llap{\pageref{l:EQS}}}\par
\noindent\hyperlink{l:ERR65}{\texttt{ERR65}}\hfill\texttt{\$2C4F~\phantom{000}\llap{\pageref{l:ERR65}}}\par
\noindent\hyperlink{l:ERR66}{\texttt{ERR66}}\hfill\texttt{\$2C49~\phantom{000}\llap{\pageref{l:ERR66}}}\par
\noindent\hyperlink{l:ERRIDX}{\texttt{ERRIDX}}\hfill\texttt{\$2D4F~\phantom{000}\llap{\pageref{l:ERRIDX}}}\par
\noindent\hyperlink{l:ERRMSG}{\texttt{ERRMSG}}\hfill\texttt{\$069B~\phantom{000}\llap{\pageref{l:ERRMSG}}}\par
\noindent\hyperlink{l:ERROR}{\texttt{ERROR}}\hfill\texttt{\$062E~\phantom{000}\llap{\pageref{l:ERROR}}}\par
\noindent\hyperlink{l:ERRSTP}{\texttt{ERRSTP}}\hfill\texttt{\$0669~\phantom{000}\llap{\pageref{l:ERRSTP}}}\par
\noindent\hyperlink{l:ERRTRP}{\texttt{ERRTRP}}\hfill\texttt{\$0659~\phantom{000}\llap{\pageref{l:ERRTRP}}}\par
\noindent\hyperlink{l:ESCARG}{\texttt{ESCARG}}\hfill\texttt{\$0227~\phantom{000}\llap{\pageref{l:ESCARG}}}\par
\noindent\hyperlink{l:EXABS}{\texttt{EXABS}}\hfill\texttt{\$2D7F~\phantom{000}\llap{\pageref{l:EXABS}}}\par
\noindent\hyperlink{l:EXADD}{\texttt{EXADD}}\hfill\texttt{\$2B9F~\phantom{000}\llap{\pageref{l:EXADD}}}\par
\noindent\hyperlink{l:EXASC}{\texttt{EXASC}}\hfill\texttt{\$2EEC~\phantom{000}\llap{\pageref{l:EXASC}}}\par
\noindent\hyperlink{l:EXATN}{\texttt{EXATN}}\hfill\texttt{\$2DAF~\phantom{000}\llap{\pageref{l:EXATN}}}\par
\noindent\hyperlink{l:EXCALL}{\texttt{EXCALL}}\hfill\texttt{\$31CC~\phantom{000}\llap{\pageref{l:EXCALL}}}\par
\noindent\hyperlink{l:EXCHR}{\texttt{EXCHR}}\hfill\texttt{\$2FC1~\phantom{000}\llap{\pageref{l:EXCHR}}}\par
\noindent\hyperlink{l:EXCOMP}{\texttt{EXCOMP}}\hfill\texttt{\$2BB3~\phantom{000}\llap{\pageref{l:EXCOMP}}}\par
\noindent\hyperlink{l:EXCOS}{\texttt{EXCOS}}\hfill\texttt{\$2DA5~\phantom{000}\llap{\pageref{l:EXCOS}}}\par
\noindent\hyperlink{l:EXCUR}{\texttt{EXCUR}}\hfill\texttt{\$313B~\phantom{000}\llap{\pageref{l:EXCUR}}}\par
\noindent\hyperlink{l:EXDIV}{\texttt{EXDIV}}\hfill\texttt{\$2BAE~\phantom{000}\llap{\pageref{l:EXDIV}}}\par
\noindent\hyperlink{l:EXDOT}{\texttt{EXDOT}}\hfill\texttt{\$22D6~\phantom{000}\llap{\pageref{l:EXDOT}}}\par
\noindent\hyperlink{l:EXEND}{\texttt{EXEND}}\hfill\texttt{\$285B~\phantom{000}\llap{\pageref{l:EXEND}}}\par
\noindent\hyperlink{l:EXERR}{\texttt{EXERR}}\hfill\texttt{\$204F~\phantom{000}\llap{\pageref{l:EXERR}}}\par
\noindent\hyperlink{l:EXESKP}{\texttt{EXESKP}}\hfill\texttt{\$0DB8~\phantom{000}\llap{\pageref{l:EXESKP}}}\par
\noindent\hyperlink{l:EXEST2}{\texttt{EXEST2}}\hfill\texttt{\$0DB3~\phantom{000}\llap{\pageref{l:EXEST2}}}\par
\noindent\hyperlink{l:EXESTM}{\texttt{EXESTM}}\hfill\texttt{\$0DA3~\phantom{000}\llap{\pageref{l:EXESTM}}}\par
\noindent\hyperlink{l:EXEXP}{\texttt{EXEXP}}\hfill\texttt{\$2842~\phantom{000}\llap{\pageref{l:EXEXP}}}\par
\noindent\hyperlink{l:EXEXP1}{\texttt{EXEXP1}}\hfill\texttt{\$2DB4~\phantom{000}\llap{\pageref{l:EXEXP1}}}\par
\noindent\hyperlink{l:EXEXPH}{\texttt{EXEXPH}}\hfill\texttt{\$2841~\phantom{000}\llap{\pageref{l:EXEXPH}}}\par
\noindent\hyperlink{l:EXFIX}{\texttt{EXFIX}}\hfill\texttt{\$2D96~\phantom{000}\llap{\pageref{l:EXFIX}}}\par
\noindent\hyperlink{l:EXFUNC}{\texttt{EXFUNC}}\hfill\texttt{\$2853~\phantom{000}\llap{\pageref{l:EXFUNC}}}\par
\noindent\hyperlink{l:EXIEC}{\texttt{EXIEC}}\hfill\texttt{\$016B~\phantom{000}\llap{\pageref{l:EXIEC}}}\par
\noindent\hyperlink{l:EXINP}{\texttt{EXINP}}\hfill\texttt{\$31C1~\phantom{000}\llap{\pageref{l:EXINP}}}\par
\noindent\hyperlink{l:EXINST}{\texttt{EXINST}}\hfill\texttt{\$3075~\phantom{000}\llap{\pageref{l:EXINST}}}\par
\noindent\hyperlink{l:EXINT}{\texttt{EXINT}}\hfill\texttt{\$2D9B~\phantom{000}\llap{\pageref{l:EXINT}}}\par
\noindent\hyperlink{l:EXLEFT}{\texttt{EXLEFT}}\hfill\texttt{\$2F07~\phantom{000}\llap{\pageref{l:EXLEFT}}}\par
\noindent\hyperlink{l:EXLEN}{\texttt{EXLEN}}\hfill\texttt{\$2EE3~\phantom{000}\llap{\pageref{l:EXLEN}}}\par
\noindent\hyperlink{l:EXLG10}{\texttt{EXLG10}}\hfill\texttt{\$2DBE~\phantom{000}\llap{\pageref{l:EXLG10}}}\par
\noindent\hyperlink{l:EXLOG}{\texttt{EXLOG}}\hfill\texttt{\$2DB9~\phantom{000}\llap{\pageref{l:EXLOG}}}\par
\noindent\hyperlink{l:EXMID}{\texttt{EXMID}}\hfill\texttt{\$2EFD~\phantom{000}\llap{\pageref{l:EXMID}}}\par
\noindent\hyperlink{l:EXMUL}{\texttt{EXMUL}}\hfill\texttt{\$2BA9~\phantom{000}\llap{\pageref{l:EXMUL}}}\par
\noindent\hyperlink{l:EXNUM}{\texttt{EXNUM}}\hfill\texttt{\$3152~\phantom{000}\llap{\pageref{l:EXNUM}}}\par
\noindent\hyperlink{l:EXPEEK}{\texttt{EXPEEK}}\hfill\texttt{\$2B7E~\phantom{000}\llap{\pageref{l:EXPEEK}}}\par
\noindent\hyperlink{l:EXPI}{\texttt{EXPI}}\hfill\texttt{\$2D6E~\phantom{000}\llap{\pageref{l:EXPI}}}\par
\noindent\hyperlink{l:EXPSER}{\texttt{EXPSER}}\hfill\texttt{\$3C7E~\phantom{000}\llap{\pageref{l:EXPSER}}}\par
\noindent\hyperlink{l:EXPTAB1}{\texttt{EXPTAB1}}\hfill\texttt{\$288B~\phantom{000}\llap{\pageref{l:EXPTAB1}}}\par
\noindent\hyperlink{l:EXPTAB2}{\texttt{EXPTAB2}}\hfill\texttt{\$291F~\phantom{000}\llap{\pageref{l:EXPTAB2}}}\par
\noindent\hyperlink{l:EXRGHT}{\texttt{EXRGHT}}\hfill\texttt{\$2F1C~\phantom{000}\llap{\pageref{l:EXRGHT}}}\par
\noindent\hyperlink{l:EXRND}{\texttt{EXRND}}\hfill\texttt{\$31E0~\phantom{000}\llap{\pageref{l:EXRND}}}\par
\noindent\hyperlink{l:EXSGN}{\texttt{EXSGN}}\hfill\texttt{\$2D86~\phantom{000}\llap{\pageref{l:EXSGN}}}\par
\noindent\hyperlink{l:EXSIN}{\texttt{EXSIN}}\hfill\texttt{\$2DA0~\phantom{000}\llap{\pageref{l:EXSIN}}}\par
\noindent\hyperlink{l:EXSPC}{\texttt{EXSPC}}\hfill\texttt{\$30F7~\phantom{000}\llap{\pageref{l:EXSPC}}}\par
\noindent\hyperlink{l:EXSQR}{\texttt{EXSQR}}\hfill\texttt{\$2DC3~\phantom{000}\llap{\pageref{l:EXSQR}}}\par
\noindent\hyperlink{l:EXSTRG}{\texttt{EXSTRG}}\hfill\texttt{\$30FB~\phantom{000}\llap{\pageref{l:EXSTRG}}}\par
\noindent\hyperlink{l:EXSUB}{\texttt{EXSUB}}\hfill\texttt{\$2BA4~\phantom{000}\llap{\pageref{l:EXSUB}}}\par
\noindent\hyperlink{l:EXSWAP}{\texttt{EXSWAP}}\hfill\texttt{\$2B79~\phantom{000}\llap{\pageref{l:EXSWAP}}}\par
\noindent\hyperlink{l:EXTAB}{\texttt{EXTAB}}\hfill\texttt{\$3126~\phantom{000}\llap{\pageref{l:EXTAB}}}\par
\noindent\hyperlink{l:EXTAN}{\texttt{EXTAN}}\hfill\texttt{\$2DAA~\phantom{000}\llap{\pageref{l:EXTAN}}}\par
\noindent\hyperlink{l:EXTJP1}{\texttt{EXTJP1}}\hfill\texttt{\$0156~\phantom{000}\llap{\pageref{l:EXTJP1}}}\par
\noindent\hyperlink{l:EXTJP2}{\texttt{EXTJP2}}\hfill\texttt{\$017F~\phantom{000}\llap{\pageref{l:EXTJP2}}}\par
\noindent\hyperlink{l:EXVAL}{\texttt{EXVAL}}\hfill\texttt{\$3172~\phantom{000}\llap{\pageref{l:EXVAL}}}\par
\noindent\hyperlink{l:F2DPI}{\texttt{F2DPI}}\hfill\texttt{\$3FA5~\phantom{000}\llap{\pageref{l:F2DPI}}}\par
\noindent\hyperlink{l:FACARY}{\texttt{FACARY}}\hfill\texttt{\$1734~\phantom{000}\llap{\pageref{l:FACARY}}}\par
\noindent\hyperlink{l:FADD}{\texttt{FADD}}\hfill\texttt{\$33F7~\phantom{000}\llap{\pageref{l:FADD}}}\par
\noindent\hyperlink{l:FADDC}{\texttt{FADDC}}\hfill\texttt{\$3402~\phantom{000}\llap{\pageref{l:FADDC}}}\par
\noindent\hyperlink{l:FADOT}{\texttt{FADOT}}\hfill\texttt{\$17BB~\phantom{000}\llap{\pageref{l:FADOT}}}\par
\noindent\hyperlink{l:FAEDGT}{\texttt{FAEDGT}}\hfill\texttt{\$17EC~\phantom{000}\llap{\pageref{l:FAEDGT}}}\par
\noindent\hyperlink{l:FAEDIG}{\texttt{FAEDIG}}\hfill\texttt{\$16E3~\phantom{000}\llap{\pageref{l:FAEDIG}}}\par
\noindent\hyperlink{l:FAEFMT}{\texttt{FAEFMT}}\hfill\texttt{\$16D1~\phantom{000}\llap{\pageref{l:FAEFMT}}}\par
\noindent\hyperlink{l:FAEND}{\texttt{FAEND}}\hfill\texttt{\$17CE~\phantom{000}\llap{\pageref{l:FAEND}}}\par
\noindent\hyperlink{l:FAEOUT}{\texttt{FAEOUT}}\hfill\texttt{\$17E3~\phantom{000}\llap{\pageref{l:FAEOUT}}}\par
\noindent\hyperlink{l:FAEPOS}{\texttt{FAEPOS}}\hfill\texttt{\$1805~\phantom{000}\llap{\pageref{l:FAEPOS}}}\par
\noindent\hyperlink{l:FAEPRN}{\texttt{FAEPRN}}\hfill\texttt{\$180D~\phantom{000}\llap{\pageref{l:FAEPRN}}}\par
\noindent\hyperlink{l:FAEXP}{\texttt{FAEXP}}\hfill\texttt{\$16C1~\phantom{000}\llap{\pageref{l:FAEXP}}}\par
\noindent\hyperlink{l:FAEXPN}{\texttt{FAEXPN}}\hfill\texttt{\$17F1~\phantom{000}\llap{\pageref{l:FAEXPN}}}\par
\noindent\hyperlink{l:FAFIX}{\texttt{FAFIX}}\hfill\texttt{\$16E7~\phantom{000}\llap{\pageref{l:FAFIX}}}\par
\noindent\hyperlink{l:FAFIX1}{\texttt{FAFIX1}}\hfill\texttt{\$16F5~\phantom{000}\llap{\pageref{l:FAFIX1}}}\par
\noindent\hyperlink{l:FAFMT}{\texttt{FAFMT}}\hfill\texttt{\$16D5~\phantom{000}\llap{\pageref{l:FAFMT}}}\par
\noindent\hyperlink{l:FAFORM}{\texttt{FAFORM}}\hfill\texttt{\$17A6~\phantom{000}\llap{\pageref{l:FAFORM}}}\par
\noindent\hyperlink{l:FAFRAC}{\texttt{FAFRAC}}\hfill\texttt{\$17C9~\phantom{000}\llap{\pageref{l:FAFRAC}}}\par
\noindent\hyperlink{l:FAINT}{\texttt{FAINT}}\hfill\texttt{\$17D4~\phantom{000}\llap{\pageref{l:FAINT}}}\par
\noindent\hyperlink{l:FALEAD}{\texttt{FALEAD}}\hfill\texttt{\$17BF~\phantom{000}\llap{\pageref{l:FALEAD}}}\par
\noindent\hyperlink{l:FALIGN}{\texttt{FALIGN}}\hfill\texttt{\$36BB~\phantom{000}\llap{\pageref{l:FALIGN}}}\par
\noindent\hyperlink{l:FAOUT}{\texttt{FAOUT}}\hfill\texttt{\$178A~\phantom{000}\llap{\pageref{l:FAOUT}}}\par
\noindent\hyperlink{l:FAOVF1}{\texttt{FAOVF1}}\hfill\texttt{\$1778~\phantom{000}\llap{\pageref{l:FAOVF1}}}\par
\noindent\hyperlink{l:FAOVF2}{\texttt{FAOVF2}}\hfill\texttt{\$1783~\phantom{000}\llap{\pageref{l:FAOVF2}}}\par
\noindent\hyperlink{l:FAOVFL}{\texttt{FAOVFL}}\hfill\texttt{\$1766~\phantom{000}\llap{\pageref{l:FAOVFL}}}\par
\noindent\hyperlink{l:FAPOS}{\texttt{FAPOS}}\hfill\texttt{\$16C9~\phantom{000}\llap{\pageref{l:FAPOS}}}\par
\noindent\hyperlink{l:FARITH}{\texttt{FARITH}}\hfill\texttt{\$2ACD~\phantom{000}\llap{\pageref{l:FARITH}}}\par
\noindent\hyperlink{l:FARND1}{\texttt{FARND1}}\hfill\texttt{\$1741~\phantom{000}\llap{\pageref{l:FARND1}}}\par
\noindent\hyperlink{l:FAROU1}{\texttt{FAROU1}}\hfill\texttt{\$1721~\phantom{000}\llap{\pageref{l:FAROU1}}}\par
\noindent\hyperlink{l:FAROUN}{\texttt{FAROUN}}\hfill\texttt{\$1719~\phantom{000}\llap{\pageref{l:FAROUN}}}\par
\noindent\hyperlink{l:FASGN1}{\texttt{FASGN1}}\hfill\texttt{\$17A4~\phantom{000}\llap{\pageref{l:FASGN1}}}\par
\noindent\hyperlink{l:FASIGN}{\texttt{FASIGN}}\hfill\texttt{\$1791~\phantom{000}\llap{\pageref{l:FASIGN}}}\par
\noindent\hyperlink{l:FASMAL}{\texttt{FASMAL}}\hfill\texttt{\$16FD~\phantom{000}\llap{\pageref{l:FASMAL}}}\par
\noindent\hyperlink{l:FASTR1}{\texttt{FASTR1}}\hfill\texttt{\$1748~\phantom{000}\llap{\pageref{l:FASTR1}}}\par
\noindent\hyperlink{l:FASTRP}{\texttt{FASTRP}}\hfill\texttt{\$1747~\phantom{000}\llap{\pageref{l:FASTRP}}}\par
\noindent\hyperlink{l:FATINY}{\texttt{FATINY}}\hfill\texttt{\$1703~\phantom{000}\llap{\pageref{l:FATINY}}}\par
\noindent\hyperlink{l:FATN}{\texttt{FATN}}\hfill\texttt{\$3B77~\phantom{000}\llap{\pageref{l:FATN}}}\par
\noindent\hyperlink{l:FAZER1}{\texttt{FAZER1}}\hfill\texttt{\$175C~\phantom{000}\llap{\pageref{l:FAZER1}}}\par
\noindent\hyperlink{l:FAZERO}{\texttt{FAZERO}}\hfill\texttt{\$1751~\phantom{000}\llap{\pageref{l:FAZERO}}}\par
\noindent\hyperlink{l:FAZERS}{\texttt{FAZERS}}\hfill\texttt{\$16C8~\phantom{000}\llap{\pageref{l:FAZERS}}}\par
\noindent\hyperlink{l:FCLR5}{\texttt{FCLR5}}\hfill\texttt{\$3C39~\phantom{000}\llap{\pageref{l:FCLR5}}}\par
\noindent\hyperlink{l:FCMP}{\texttt{FCMP}}\hfill\texttt{\$2A93~\phantom{000}\llap{\pageref{l:FCMP}}}\par
\noindent\hyperlink{l:FCMPST}{\texttt{FCMPST}}\hfill\texttt{\$2A7A~\phantom{000}\llap{\pageref{l:FCMPST}}}\par
\noindent\hyperlink{l:FCOPY}{\texttt{FCOPY}}\hfill\texttt{\$36FC~\phantom{000}\llap{\pageref{l:FCOPY}}}\par
\noindent\hyperlink{l:FCOS}{\texttt{FCOS}}\hfill\texttt{\$3D70~\phantom{000}\llap{\pageref{l:FCOS}}}\par
\noindent\hyperlink{l:FDIV}{\texttt{FDIV}}\hfill\texttt{\$35EE~\phantom{000}\llap{\pageref{l:FDIV}}}\par
\noindent\hyperlink{l:FDLOOP}{\texttt{FDLOOP}}\hfill\texttt{\$3622~\phantom{000}\llap{\pageref{l:FDLOOP}}}\par
\noindent\hyperlink{l:FDRND}{\texttt{FDRND}}\hfill\texttt{\$3689~\phantom{000}\llap{\pageref{l:FDRND}}}\par
\noindent\hyperlink{l:FENTER}{\texttt{FENTER}}\hfill\texttt{\$3766~\phantom{000}\llap{\pageref{l:FENTER}}}\par
\noindent\hyperlink{l:FEXP}{\texttt{FEXP}}\hfill\texttt{\$3C03~\phantom{000}\llap{\pageref{l:FEXP}}}\par
\noindent\hyperlink{l:FFIX}{\texttt{FFIX}}\hfill\texttt{\$2DD5~\phantom{000}\llap{\pageref{l:FFIX}}}\par
\noindent\hyperlink{l:FFUNC}{\texttt{FFUNC}}\hfill\texttt{\$2DC6~\phantom{000}\llap{\pageref{l:FFUNC}}}\par
\noindent\hyperlink{l:FINDLN}{\texttt{FINDLN}}\hfill\texttt{\$0F39~\phantom{000}\llap{\pageref{l:FINDLN}}}\par
\noindent\hyperlink{l:FINT}{\texttt{FINT}}\hfill\texttt{\$3C7F~\phantom{000}\llap{\pageref{l:FINT}}}\par
\noindent\hyperlink{l:FINT1}{\texttt{FINT1}}\hfill\texttt{\$3C8B~\phantom{000}\llap{\pageref{l:FINT1}}}\par
\noindent\hyperlink{l:FINT2}{\texttt{FINT2}}\hfill\texttt{\$3C8E~\phantom{000}\llap{\pageref{l:FINT2}}}\par
\noindent\hyperlink{l:FINTFL}{\texttt{FINTFL}}\hfill\texttt{\$3CA7~\phantom{000}\llap{\pageref{l:FINTFL}}}\par
\noindent\hyperlink{l:FINTLP}{\texttt{FINTLP}}\hfill\texttt{\$3C93~\phantom{000}\llap{\pageref{l:FINTLP}}}\par
\noindent\hyperlink{l:FLDX}{\texttt{FLDX}}\hfill\texttt{\$3BC5~\phantom{000}\llap{\pageref{l:FLDX}}}\par
\noindent\hyperlink{l:FLGE}{\texttt{FLGE}}\hfill\texttt{\$3D31~\phantom{000}\llap{\pageref{l:FLGE}}}\par
\noindent\hyperlink{l:FLN10}{\texttt{FLN10}}\hfill\texttt{\$3D36~\phantom{000}\llap{\pageref{l:FLN10}}}\par
\noindent\hyperlink{l:FLOG}{\texttt{FLOG}}\hfill\texttt{\$3CC4~\phantom{000}\llap{\pageref{l:FLOG}}}\par
\noindent\hyperlink{l:FLOG10}{\texttt{FLOG10}}\hfill\texttt{\$3CB7~\phantom{000}\llap{\pageref{l:FLOG10}}}\par
\noindent\hyperlink{l:FLTAS6}{\texttt{FLTAS6}}\hfill\texttt{\$1677~\phantom{000}\llap{\pageref{l:FLTAS6}}}\par
\noindent\hyperlink{l:FLTASC}{\texttt{FLTASC}}\hfill\texttt{\$1679~\phantom{000}\llap{\pageref{l:FLTASC}}}\par
\noindent\hyperlink{l:FLTCH}{\texttt{FLTCH}}\hfill\texttt{\$15C4~\phantom{000}\llap{\pageref{l:FLTCH}}}\par
\noindent\hyperlink{l:FLTEXP}{\texttt{FLTEXP}}\hfill\texttt{\$1642~\phantom{000}\llap{\pageref{l:FLTEXP}}}\par
\noindent\hyperlink{l:FLTSAV}{\texttt{FLTSAV}}\hfill\texttt{\$1670~\phantom{000}\llap{\pageref{l:FLTSAV}}}\par
\noindent\hyperlink{l:FMLN2}{\texttt{FMLN2}}\hfill\texttt{\$3D3B~\phantom{000}\llap{\pageref{l:FMLN2}}}\par
\noindent\hyperlink{l:FMLOOP}{\texttt{FMLOOP}}\hfill\texttt{\$356C~\phantom{000}\llap{\pageref{l:FMLOOP}}}\par
\noindent\hyperlink{l:FMRND}{\texttt{FMRND}}\hfill\texttt{\$359D~\phantom{000}\llap{\pageref{l:FMRND}}}\par
\noindent\hyperlink{l:FMUL}{\texttt{FMUL}}\hfill\texttt{\$34DF~\phantom{000}\llap{\pageref{l:FMUL}}}\par
\noindent\hyperlink{l:FMULE}{\texttt{FMULE}}\hfill\texttt{\$34FA~\phantom{000}\llap{\pageref{l:FMULE}}}\par
\noindent\hyperlink{l:FN}{\texttt{FN}}\hfill\texttt{\$1B2A~\phantom{000}\llap{\pageref{l:FN}}}\par
\noindent\hyperlink{l:FNAME}{\texttt{FNAME}}\hfill\texttt{\$1478~\phantom{000}\llap{\pageref{l:FNAME}}}\par
\noindent\hyperlink{l:FNDFIL}{\texttt{FNDFIL}}\hfill\texttt{\$20E1~\phantom{000}\llap{\pageref{l:FNDFIL}}}\par
\noindent\hyperlink{l:FNEG}{\texttt{FNEG}}\hfill\texttt{\$3BC0~\phantom{000}\llap{\pageref{l:FNEG}}}\par
\noindent\hyperlink{l:FNEGX}{\texttt{FNEGX}}\hfill\texttt{\$3BBD~\phantom{000}\llap{\pageref{l:FNEGX}}}\par
\noindent\hyperlink{l:FNKW}{\texttt{FNKW}}\hfill\texttt{\$2431~\phantom{000}\llap{\pageref{l:FNKW}}}\par
\noindent\hyperlink{l:FNKWTX}{\texttt{FNKWTX}}\hfill\texttt{\$13F1~\phantom{000}\llap{\pageref{l:FNKWTX}}}\par
\noindent\hyperlink{l:FNORM}{\texttt{FNORM}}\hfill\texttt{\$3494~\phantom{000}\llap{\pageref{l:FNORM}}}\par
\noindent\hyperlink{l:FNPAR}{\texttt{FNPAR}}\hfill\texttt{\$33AB~\phantom{000}\llap{\pageref{l:FNPAR}}}\par
\noindent\hyperlink{l:FNPART}{\texttt{FNPART}}\hfill\texttt{\$14DA~\phantom{000}\llap{\pageref{l:FNPART}}}\par
\noindent\hyperlink{l:FNRES}{\texttt{FNRES}}\hfill\texttt{\$1B7B~\phantom{000}\llap{\pageref{l:FNRES}}}\par
\noindent\hyperlink{l:FNTEXT}{\texttt{FNTEXT}}\hfill\texttt{\$13E9~\phantom{000}\llap{\pageref{l:FNTEXT}}}\par
\noindent\hyperlink{l:FNTYPE}{\texttt{FNTYPE}}\hfill\texttt{\$24F2~\phantom{000}\llap{\pageref{l:FNTYPE}}}\par
\noindent\hyperlink{l:FONE}{\texttt{FONE}}\hfill\texttt{\$3FBB~\phantom{000}\llap{\pageref{l:FONE}}}\par
\noindent\hyperlink{l:FOR}{\texttt{FOR}}\hfill\texttt{\$1CC2~\phantom{000}\llap{\pageref{l:FOR}}}\par
\noindent\hyperlink{l:FORTSF}{\texttt{FORTSF}}\hfill\texttt{\$1D88~\phantom{000}\llap{\pageref{l:FORTSF}}}\par
\noindent\hyperlink{l:FORTSI}{\texttt{FORTSI}}\hfill\texttt{\$1D59~\phantom{000}\llap{\pageref{l:FORTSI}}}\par
\noindent\hyperlink{l:FOVFL}{\texttt{FOVFL}}\hfill\texttt{\$3758~\phantom{000}\llap{\pageref{l:FOVFL}}}\par
\noindent\hyperlink{l:FPI}{\texttt{FPI}}\hfill\texttt{\$3FA6~\phantom{000}\llap{\pageref{l:FPI}}}\par
\noindent\hyperlink{l:FPIH}{\texttt{FPIH}}\hfill\texttt{\$3FA0~\phantom{000}\llap{\pageref{l:FPIH}}}\par
\noindent\hyperlink{l:FPIQ}{\texttt{FPIQ}}\hfill\texttt{\$3FB4~\phantom{000}\llap{\pageref{l:FPIQ}}}\par
\noindent\hyperlink{l:FPOLY}{\texttt{FPOLY}}\hfill\texttt{\$3F61~\phantom{000}\llap{\pageref{l:FPOLY}}}\par
\noindent\hyperlink{l:FPOW}{\texttt{FPOW}}\hfill\texttt{\$2E07~\phantom{000}\llap{\pageref{l:FPOW}}}\par
\noindent\hyperlink{l:FSIN}{\texttt{FSIN}}\hfill\texttt{\$3D77~\phantom{000}\llap{\pageref{l:FSIN}}}\par
\noindent\hyperlink{l:FSINC}{\texttt{FSINC}}\hfill\texttt{\$3D78~\phantom{000}\llap{\pageref{l:FSINC}}}\par
\noindent\hyperlink{l:FSQR}{\texttt{FSQR}}\hfill\texttt{\$3DF5~\phantom{000}\llap{\pageref{l:FSQR}}}\par
\noindent\hyperlink{l:FSTOR1}{\texttt{FSTOR1}}\hfill\texttt{\$3437~\phantom{000}\llap{\pageref{l:FSTOR1}}}\par
\noindent\hyperlink{l:FSTORE}{\texttt{FSTORE}}\hfill\texttt{\$3434~\phantom{000}\llap{\pageref{l:FSTORE}}}\par
\noindent\hyperlink{l:FSUB}{\texttt{FSUB}}\hfill\texttt{\$3460~\phantom{000}\llap{\pageref{l:FSUB}}}\par
\noindent\hyperlink{l:FSUBM}{\texttt{FSUBM}}\hfill\texttt{\$3473~\phantom{000}\llap{\pageref{l:FSUBM}}}\par
\noindent\hyperlink{l:FTAN}{\texttt{FTAN}}\hfill\texttt{\$3EAD~\phantom{000}\llap{\pageref{l:FTAN}}}\par
\noindent\hyperlink{l:FTIERR}{\texttt{FTIERR}}\hfill\texttt{\$3EF1~\phantom{000}\llap{\pageref{l:FTIERR}}}\par
\noindent\hyperlink{l:FTILP}{\texttt{FTILP}}\hfill\texttt{\$3EDE~\phantom{000}\llap{\pageref{l:FTILP}}}\par
\noindent\hyperlink{l:FTISGN}{\texttt{FTISGN}}\hfill\texttt{\$3EF8~\phantom{000}\llap{\pageref{l:FTISGN}}}\par
\noindent\hyperlink{l:FTOI}{\texttt{FTOI}}\hfill\texttt{\$2DE2~\phantom{000}\llap{\pageref{l:FTOI}}}\par
\noindent\hyperlink{l:FTOINT}{\texttt{FTOINT}}\hfill\texttt{\$3ECA~\phantom{000}\llap{\pageref{l:FTOINT}}}\par
\noindent\hyperlink{l:FZERO}{\texttt{FZERO}}\hfill\texttt{\$375B~\phantom{000}\llap{\pageref{l:FZERO}}}\par
\noindent\hyperlink{l:GEF}{\texttt{GEF}}\hfill\texttt{\$2A75~\phantom{000}\llap{\pageref{l:GEF}}}\par
\noindent\hyperlink{l:GEI}{\texttt{GEI}}\hfill\texttt{\$2B4A~\phantom{000}\llap{\pageref{l:GEI}}}\par
\noindent\hyperlink{l:GES}{\texttt{GES}}\hfill\texttt{\$300D~\phantom{000}\llap{\pageref{l:GES}}}\par
\noindent\hyperlink{l:GET}{\texttt{GET}}\hfill\texttt{\$1DB1~\phantom{000}\llap{\pageref{l:GET}}}\par
\noindent\hyperlink{l:GETDIG}{\texttt{GETDIG}}\hfill\texttt{\$1837~\phantom{000}\llap{\pageref{l:GETDIG}}}\par
\noindent\hyperlink{l:GETEND}{\texttt{GETEND}}\hfill\texttt{\$1833~\phantom{000}\llap{\pageref{l:GETEND}}}\par
\noindent\hyperlink{l:GETF}{\texttt{GETF}}\hfill\texttt{\$29A1~\phantom{000}\llap{\pageref{l:GETF}}}\par
\noindent\hyperlink{l:GETFLT}{\texttt{GETFLT}}\hfill\texttt{\$1513~\phantom{000}\llap{\pageref{l:GETFLT}}}\par
\noindent\hyperlink{l:GETI}{\texttt{GETI}}\hfill\texttt{\$29BD~\phantom{000}\llap{\pageref{l:GETI}}}\par
\noindent\hyperlink{l:GETKEY}{\texttt{GETKEY}}\hfill\texttt{\$02EA~\phantom{000}\llap{\pageref{l:GETKEY}}}\par
\noindent\hyperlink{l:GETLN}{\texttt{GETLN}}\hfill\texttt{\$181E~\phantom{000}\llap{\pageref{l:GETLN}}}\par
\noindent\hyperlink{l:GETLN1}{\texttt{GETLN1}}\hfill\texttt{\$181D~\phantom{000}\llap{\pageref{l:GETLN1}}}\par
\noindent\hyperlink{l:GETLP}{\texttt{GETLP}}\hfill\texttt{\$182A~\phantom{000}\llap{\pageref{l:GETLP}}}\par
\noindent\hyperlink{l:GETOVF}{\texttt{GETOVF}}\hfill\texttt{\$183F~\phantom{000}\llap{\pageref{l:GETOVF}}}\par
\noindent\hyperlink{l:GETPOS}{\texttt{GETPOS}}\hfill\texttt{\$08EF~\phantom{000}\llap{\pageref{l:GETPOS}}}\par
\noindent\hyperlink{l:GETS}{\texttt{GETS}}\hfill\texttt{\$2A22~\phantom{000}\llap{\pageref{l:GETS}}}\par
\noindent\hyperlink{l:GETSFL}{\texttt{GETSFL}}\hfill\texttt{\$14FB~\phantom{000}\llap{\pageref{l:GETSFL}}}\par
\noindent\hyperlink{l:GETSIN}{\texttt{GETSIN}}\hfill\texttt{\$1841~\phantom{000}\llap{\pageref{l:GETSIN}}}\par
\noindent\hyperlink{l:GOSUB}{\texttt{GOSUB}}\hfill\texttt{\$1DE9~\phantom{000}\llap{\pageref{l:GOSUB}}}\par
\noindent\hyperlink{l:GOTO}{\texttt{GOTO}}\hfill\texttt{\$1DDB~\phantom{000}\llap{\pageref{l:GOTO}}}\par
\noindent\hyperlink{l:GTF}{\texttt{GTF}}\hfill\texttt{\$2A5E~\phantom{000}\llap{\pageref{l:GTF}}}\par
\noindent\hyperlink{l:GTI}{\texttt{GTI}}\hfill\texttt{\$2B5A~\phantom{000}\llap{\pageref{l:GTI}}}\par
\noindent\hyperlink{l:GTS}{\texttt{GTS}}\hfill\texttt{\$2FF6~\phantom{000}\llap{\pageref{l:GTS}}}\par
\noindent\hyperlink{l:IDIV}{\texttt{IDIV}}\hfill\texttt{\$3788~\phantom{000}\llap{\pageref{l:IDIV}}}\par
\noindent\hyperlink{l:IF}{\texttt{IF}}\hfill\texttt{\$1E4B~\phantom{000}\llap{\pageref{l:IF}}}\par
\noindent\hyperlink{l:IM2VEC}{\texttt{IM2VEC}}\hfill\texttt{\$0034~\phantom{000}\llap{\pageref{l:IM2VEC}}}\par
\noindent\hyperlink{l:IMUL}{\texttt{IMUL}}\hfill\texttt{\$3774~\phantom{000}\llap{\pageref{l:IMUL}}}\par
\noindent\hyperlink{l:IMULDV}{\texttt{IMULDV}}\hfill\texttt{\$2B6B~\phantom{000}\llap{\pageref{l:IMULDV}}}\par
\noindent\hyperlink{l:INIT}{\texttt{INIT}}\hfill\texttt{\$0068~\phantom{000}\llap{\pageref{l:INIT}}}\par
\noindent\hyperlink{l:INLIN2}{\texttt{INLIN2}}\hfill\texttt{\$02AD~\phantom{000}\llap{\pageref{l:INLIN2}}}\par
\noindent\hyperlink{l:INLINE}{\texttt{INLINE}}\hfill\texttt{\$02AA~\phantom{000}\llap{\pageref{l:INLINE}}}\par
\noindent\hyperlink{l:INPLIN}{\texttt{INPLIN}}\hfill\texttt{\$1F82~\phantom{000}\llap{\pageref{l:INPLIN}}}\par
\noindent\hyperlink{l:INPT}{\texttt{INPT}}\hfill\texttt{\$1E78~\phantom{000}\llap{\pageref{l:INPT}}}\par
\noindent\hyperlink{l:INPUT}{\texttt{INPUT}}\hfill\texttt{\$08DD~\phantom{000}\llap{\pageref{l:INPUT}}}\par
\noindent\hyperlink{l:INSCNV}{\texttt{INSCNV}}\hfill\texttt{\$27CF~\phantom{000}\llap{\pageref{l:INSCNV}}}\par
\noindent\hyperlink{l:INSLIN}{\texttt{INSLIN}}\hfill\texttt{\$0750~\phantom{000}\llap{\pageref{l:INSLIN}}}\par
\noindent\hyperlink{l:INTDHI}{\texttt{INTDHI}}\hfill\texttt{\$3F5C~\phantom{000}\llap{\pageref{l:INTDHI}}}\par
\noindent\hyperlink{l:INTDIG}{\texttt{INTDIG}}\hfill\texttt{\$3F3F~\phantom{000}\llap{\pageref{l:INTDIG}}}\par
\noindent\hyperlink{l:INTTOF}{\texttt{INTTOF}}\hfill\texttt{\$3F18~\phantom{000}\llap{\pageref{l:INTTOF}}}\par
\noindent\hyperlink{l:ITFPOS}{\texttt{ITFPOS}}\hfill\texttt{\$3F36~\phantom{000}\llap{\pageref{l:ITFPOS}}}\par
\noindent\hyperlink{l:ITOF}{\texttt{ITOF}}\hfill\texttt{\$2DF0~\phantom{000}\llap{\pageref{l:ITOF}}}\par
\noindent\hyperlink{l:IXOFF}{\texttt{IXOFF}}\hfill\texttt{\$375F~\phantom{000}\llap{\pageref{l:IXOFF}}}\par
\noindent\hyperlink{l:JOINA}{\texttt{JOINA}}\hfill\texttt{\$140C~\phantom{000}\llap{\pageref{l:JOINA}}}\par
\noindent\hyperlink{l:JOINC}{\texttt{JOINC}}\hfill\texttt{\$140A~\phantom{000}\llap{\pageref{l:JOINC}}}\par
\noindent\hyperlink{l:JPHL}{\texttt{JPHL}}\hfill\texttt{\$2852~\phantom{000}\llap{\pageref{l:JPHL}}}\par
\noindent\hyperlink{l:JPIX}{\texttt{JPIX}}\hfill\texttt{\$2A91~\phantom{000}\llap{\pageref{l:JPIX}}}\par
\noindent\hyperlink{l:KEYINT}{\texttt{KEYINT}}\hfill\texttt{\$031E~\phantom{000}\llap{\pageref{l:KEYINT}}}\par
\noindent\hyperlink{l:KWADD}{\texttt{KWADD}}\hfill\texttt{\$2421~\phantom{000}\llap{\pageref{l:KWADD}}}\par
\noindent\hyperlink{l:KWAND}{\texttt{KWAND}}\hfill\texttt{\$2407~\phantom{000}\llap{\pageref{l:KWAND}}}\par
\noindent\hyperlink{l:KWAS}{\texttt{KWAS}}\hfill\texttt{\$1998~\phantom{000}\llap{\pageref{l:KWAS}}}\par
\noindent\hyperlink{l:KWASF}{\texttt{KWASF}}\hfill\texttt{\$1990~\phantom{000}\llap{\pageref{l:KWASF}}}\par
\noindent\hyperlink{l:KWCMP}{\texttt{KWCMP}}\hfill\texttt{\$2411~\phantom{000}\llap{\pageref{l:KWCMP}}}\par
\noindent\hyperlink{l:KWDEFN}{\texttt{KWDEFN}}\hfill\texttt{\$196D~\phantom{000}\llap{\pageref{l:KWDEFN}}}\par
\noindent\hyperlink{l:KWELSE}{\texttt{KWELSE}}\hfill\texttt{\$1980~\phantom{000}\llap{\pageref{l:KWELSE}}}\par
\noindent\hyperlink{l:KWEQV}{\texttt{KWEQV}}\hfill\texttt{\$23F5~\phantom{000}\llap{\pageref{l:KWEQV}}}\par
\noindent\hyperlink{l:KWIMP}{\texttt{KWIMP}}\hfill\texttt{\$23FA~\phantom{000}\llap{\pageref{l:KWIMP}}}\par
\noindent\hyperlink{l:KWLIST}{\texttt{KWLIST}}\hfill\texttt{\$1973~\phantom{000}\llap{\pageref{l:KWLIST}}}\par
\noindent\hyperlink{l:KWMUL}{\texttt{KWMUL}}\hfill\texttt{\$2426~\phantom{000}\llap{\pageref{l:KWMUL}}}\par
\noindent\hyperlink{l:KWNOT}{\texttt{KWNOT}}\hfill\texttt{\$240C~\phantom{000}\llap{\pageref{l:KWNOT}}}\par
\noindent\hyperlink{l:KWOR}{\texttt{KWOR}}\hfill\texttt{\$23FF~\phantom{000}\llap{\pageref{l:KWOR}}}\par
\noindent\hyperlink{l:KWPOW}{\texttt{KWPOW}}\hfill\texttt{\$242B~\phantom{000}\llap{\pageref{l:KWPOW}}}\par
\noindent\hyperlink{l:KWSRCH}{\texttt{KWSRCH}}\hfill\texttt{\$09F5~\phantom{000}\llap{\pageref{l:KWSRCH}}}\par
\noindent\hyperlink{l:KWSTEP}{\texttt{KWSTEP}}\hfill\texttt{\$1986~\phantom{000}\llap{\pageref{l:KWSTEP}}}\par
\noindent\hyperlink{l:KWTHEN}{\texttt{KWTHEN}}\hfill\texttt{\$197A~\phantom{000}\llap{\pageref{l:KWTHEN}}}\par
\noindent\hyperlink{l:KWTO}{\texttt{KWTO}}\hfill\texttt{\$198C~\phantom{000}\llap{\pageref{l:KWTO}}}\par
\noindent\hyperlink{l:LEF}{\texttt{LEF}}\hfill\texttt{\$2A67~\phantom{000}\llap{\pageref{l:LEF}}}\par
\noindent\hyperlink{l:LEI}{\texttt{LEI}}\hfill\texttt{\$2B56~\phantom{000}\llap{\pageref{l:LEI}}}\par
\noindent\hyperlink{l:LES}{\texttt{LES}}\hfill\texttt{\$2FFF~\phantom{000}\llap{\pageref{l:LES}}}\par
\noindent\hyperlink{l:LET}{\texttt{LET}}\hfill\texttt{\$1BFD~\phantom{000}\llap{\pageref{l:LET}}}\par
\noindent\hyperlink{l:LINFD}{\texttt{LINFD}}\hfill\texttt{\$0238~\phantom{000}\llap{\pageref{l:LINFD}}}\par
\noindent\hyperlink{l:LNKADR}{\texttt{LNKADR}}\hfill\texttt{\$10A4~\phantom{000}\llap{\pageref{l:LNKADR}}}\par
\noindent\hyperlink{l:LNKDEF}{\texttt{LNKDEF}}\hfill\texttt{\$0F59~\phantom{000}\llap{\pageref{l:LNKDEF}}}\par
\noindent\hyperlink{l:LNKELM}{\texttt{LNKELM}}\hfill\texttt{\$0EAA~\phantom{000}\llap{\pageref{l:LNKELM}}}\par
\noindent\hyperlink{l:LNKEXP}{\texttt{LNKEXP}}\hfill\texttt{\$1032~\phantom{000}\llap{\pageref{l:LNKEXP}}}\par
\noindent\hyperlink{l:LNKFN}{\texttt{LNKFN}}\hfill\texttt{\$111D~\phantom{000}\llap{\pageref{l:LNKFN}}}\par
\noindent\hyperlink{l:LNKFOR}{\texttt{LNKFOR}}\hfill\texttt{\$0F7C~\phantom{000}\llap{\pageref{l:LNKFOR}}}\par
\noindent\hyperlink{l:LNKLST}{\texttt{LNKLST}}\hfill\texttt{\$106B~\phantom{000}\llap{\pageref{l:LNKLST}}}\par
\noindent\hyperlink{l:LNKNEW}{\texttt{LNKNEW}}\hfill\texttt{\$108E~\phantom{000}\llap{\pageref{l:LNKNEW}}}\par
\noindent\hyperlink{l:LNKNXT}{\texttt{LNKNXT}}\hfill\texttt{\$0FDB~\phantom{000}\llap{\pageref{l:LNKNXT}}}\par
\noindent\hyperlink{l:LNKPAR}{\texttt{LNKPAR}}\hfill\texttt{\$1050~\phantom{000}\llap{\pageref{l:LNKPAR}}}\par
\noindent\hyperlink{l:LNKREF}{\texttt{LNKREF}}\hfill\texttt{\$0D01~\phantom{000}\llap{\pageref{l:LNKREF}}}\par
\noindent\hyperlink{l:LNKSKP}{\texttt{LNKSKP}}\hfill\texttt{\$102E~\phantom{000}\llap{\pageref{l:LNKSKP}}}\par
\noindent\hyperlink{l:LNKSRC}{\texttt{LNKSRC}}\hfill\texttt{\$1074~\phantom{000}\llap{\pageref{l:LNKSRC}}}\par
\noindent\hyperlink{l:LNKVAR}{\texttt{LNKVAR}}\hfill\texttt{\$103C~\phantom{000}\llap{\pageref{l:LNKVAR}}}\par
\noindent\hyperlink{l:LNKWR}{\texttt{LNKWR}}\hfill\texttt{\$10B0~\phantom{000}\llap{\pageref{l:LNKWR}}}\par
\noindent\hyperlink{l:LOGDBL}{\texttt{LOGDBL}}\hfill\texttt{\$3CCE~\phantom{000}\llap{\pageref{l:LOGDBL}}}\par
\noindent\hyperlink{l:LOGEXP}{\texttt{LOGEXP}}\hfill\texttt{\$3D0D~\phantom{000}\llap{\pageref{l:LOGEXP}}}\par
\noindent\hyperlink{l:LOGK}{\texttt{LOGK}}\hfill\texttt{\$3D03~\phantom{000}\llap{\pageref{l:LOGK}}}\par
\noindent\hyperlink{l:LOGRED}{\texttt{LOGRED}}\hfill\texttt{\$3CE3~\phantom{000}\llap{\pageref{l:LOGRED}}}\par
\noindent\hyperlink{l:LOGSER}{\texttt{LOGSER}}\hfill\texttt{\$3D6F~\phantom{000}\llap{\pageref{l:LOGSER}}}\par
\noindent\hyperlink{l:LTF}{\texttt{LTF}}\hfill\texttt{\$2A70~\phantom{000}\llap{\pageref{l:LTF}}}\par
\noindent\hyperlink{l:LTI}{\texttt{LTI}}\hfill\texttt{\$2B3E~\phantom{000}\llap{\pageref{l:LTI}}}\par
\noindent\hyperlink{l:LTS}{\texttt{LTS}}\hfill\texttt{\$3008~\phantom{000}\llap{\pageref{l:LTS}}}\par
\noindent\hyperlink{l:MADD3}{\texttt{MADD3}}\hfill\texttt{\$3447~\phantom{000}\llap{\pageref{l:MADD3}}}\par
\noindent\hyperlink{l:MEMFUL}{\texttt{MEMFUL}}\hfill\texttt{\$3398~\phantom{000}\llap{\pageref{l:MEMFUL}}}\par
\noindent\hyperlink{l:MEMRES}{\texttt{MEMRES}}\hfill\texttt{\$335A~\phantom{000}\llap{\pageref{l:MEMRES}}}\par
\noindent\hyperlink{l:MSHL}{\texttt{MSHL}}\hfill\texttt{\$3736~\phantom{000}\llap{\pageref{l:MSHL}}}\par
\noindent\hyperlink{l:MSHR}{\texttt{MSHR}}\hfill\texttt{\$3714~\phantom{000}\llap{\pageref{l:MSHR}}}\par
\noindent\hyperlink{l:MSUB3}{\texttt{MSUB3}}\hfill\texttt{\$34CB~\phantom{000}\llap{\pageref{l:MSUB3}}}\par
\noindent\hyperlink{l:MUL10A}{\texttt{MUL10A}}\hfill\texttt{\$3F00~\phantom{000}\llap{\pageref{l:MUL10A}}}\par
\noindent\hyperlink{l:MUL10X}{\texttt{MUL10X}}\hfill\texttt{\$3F16~\phantom{000}\llap{\pageref{l:MUL10X}}}\par
\noindent\hyperlink{l:MULF}{\texttt{MULF}}\hfill\texttt{\$2AC3~\phantom{000}\llap{\pageref{l:MULF}}}\par
\noindent\hyperlink{l:MULI}{\texttt{MULI}}\hfill\texttt{\$2B67~\phantom{000}\llap{\pageref{l:MULI}}}\par
\noindent\hyperlink{l:NEF}{\texttt{NEF}}\hfill\texttt{\$2A55~\phantom{000}\llap{\pageref{l:NEF}}}\par
\noindent\hyperlink{l:NEGF}{\texttt{NEGF}}\hfill\texttt{\$2A43~\phantom{000}\llap{\pageref{l:NEGF}}}\par
\noindent\hyperlink{l:NEGHL}{\texttt{NEGHL}}\hfill\texttt{\$2AEE~\phantom{000}\llap{\pageref{l:NEGHL}}}\par
\noindent\hyperlink{l:NEGI}{\texttt{NEGI}}\hfill\texttt{\$2AE9~\phantom{000}\llap{\pageref{l:NEGI}}}\par
\noindent\hyperlink{l:NEI}{\texttt{NEI}}\hfill\texttt{\$2B34~\phantom{000}\llap{\pageref{l:NEI}}}\par
\noindent\hyperlink{l:NES}{\texttt{NES}}\hfill\texttt{\$2FED~\phantom{000}\llap{\pageref{l:NES}}}\par
\noindent\hyperlink{l:NEW}{\texttt{NEW}}\hfill\texttt{\$0A70~\phantom{000}\llap{\pageref{l:NEW}}}\par
\noindent\hyperlink{l:NEWLIN}{\texttt{NEWLIN}}\hfill\texttt{\$0742~\phantom{000}\llap{\pageref{l:NEWLIN}}}\par
\noindent\hyperlink{l:NEXT}{\texttt{NEXT}}\hfill\texttt{\$1D21~\phantom{000}\llap{\pageref{l:NEXT}}}\par
\noindent\hyperlink{l:NMI}{\texttt{NMI}}\hfill\texttt{\$0066~\phantom{000}\llap{\pageref{l:NMI}}}\par
\noindent\hyperlink{l:NMICLK}{\texttt{NMICLK}}\hfill\texttt{\$0047~\phantom{000}\llap{\pageref{l:NMICLK}}}\par
\noindent\hyperlink{l:NOCODE}{\texttt{NOCODE}}\hfill\texttt{\$22B5~\phantom{000}\llap{\pageref{l:NOCODE}}}\par
\noindent\hyperlink{l:NODATA}{\texttt{NODATA}}\hfill\texttt{\$10FB~\phantom{000}\llap{\pageref{l:NODATA}}}\par
\noindent\hyperlink{l:NODIR}{\texttt{NODIR}}\hfill\texttt{\$0DBF~\phantom{000}\llap{\pageref{l:NODIR}}}\par
\noindent\hyperlink{l:NOTRAC}{\texttt{NOTRAC}}\hfill\texttt{\$232A~\phantom{000}\llap{\pageref{l:NOTRAC}}}\par
\noindent\hyperlink{l:NXTLEV}{\texttt{NXTLEV}}\hfill\texttt{\$27E8~\phantom{000}\llap{\pageref{l:NXTLEV}}}\par
\noindent\hyperlink{l:ON}{\texttt{ON}}\hfill\texttt{\$2009~\phantom{000}\llap{\pageref{l:ON}}}\par
\noindent\hyperlink{l:ONERRG}{\texttt{ONERRG}}\hfill\texttt{\$2041~\phantom{000}\llap{\pageref{l:ONERRG}}}\par
\noindent\hyperlink{l:OPAND}{\texttt{OPAND}}\hfill\texttt{\$2B14~\phantom{000}\llap{\pageref{l:OPAND}}}\par
\noindent\hyperlink{l:OPEN}{\texttt{OPEN}}\hfill\texttt{\$07E2~\phantom{000}\llap{\pageref{l:OPEN}}}\par
\noindent\hyperlink{l:OPENST}{\texttt{OPENST}}\hfill\texttt{\$2074~\phantom{000}\llap{\pageref{l:OPENST}}}\par
\noindent\hyperlink{l:OPEQV}{\texttt{OPEQV}}\hfill\texttt{\$2AF5~\phantom{000}\llap{\pageref{l:OPEQV}}}\par
\noindent\hyperlink{l:OPIMP}{\texttt{OPIMP}}\hfill\texttt{\$2B00~\phantom{000}\llap{\pageref{l:OPIMP}}}\par
\noindent\hyperlink{l:OPNOT}{\texttt{OPNOT}}\hfill\texttt{\$2B1D~\phantom{000}\llap{\pageref{l:OPNOT}}}\par
\noindent\hyperlink{l:OPNSTR}{\texttt{OPNSTR}}\hfill\texttt{\$07C2~\phantom{000}\llap{\pageref{l:OPNSTR}}}\par
\noindent\hyperlink{l:OPOR}{\texttt{OPOR}}\hfill\texttt{\$2B0B~\phantom{000}\llap{\pageref{l:OPOR}}}\par
\noindent\hyperlink{l:OPTEXT}{\texttt{OPTEXT}}\hfill\texttt{\$13D6~\phantom{000}\llap{\pageref{l:OPTEXT}}}\par
\noindent\hyperlink{l:OPTYPE}{\texttt{OPTYPE}}\hfill\texttt{\$27A0~\phantom{000}\llap{\pageref{l:OPTYPE}}}\par
\noindent\hyperlink{l:OPXOR}{\texttt{OPXOR}}\hfill\texttt{\$2B21~\phantom{000}\llap{\pageref{l:OPXOR}}}\par
\noindent\hyperlink{l:OUTST}{\texttt{OUTST}}\hfill\texttt{\$212A~\phantom{000}\llap{\pageref{l:OUTST}}}\par
\noindent\hyperlink{l:POKE}{\texttt{POKE}}\hfill\texttt{\$2144~\phantom{000}\llap{\pageref{l:POKE}}}\par
\noindent\hyperlink{l:POLYLP}{\texttt{POLYLP}}\hfill\texttt{\$3F7D~\phantom{000}\llap{\pageref{l:POLYLP}}}\par
\noindent\hyperlink{l:POWF}{\texttt{POWF}}\hfill\texttt{\$2E00~\phantom{000}\llap{\pageref{l:POWF}}}\par
\noindent\hyperlink{l:POWFI}{\texttt{POWFI}}\hfill\texttt{\$2E5E~\phantom{000}\llap{\pageref{l:POWFI}}}\par
\noindent\hyperlink{l:POWI}{\texttt{POWI}}\hfill\texttt{\$2E18~\phantom{000}\llap{\pageref{l:POWI}}}\par
\noindent\hyperlink{l:PRCRLF}{\texttt{PRCRLF}}\hfill\texttt{\$21F6~\phantom{000}\llap{\pageref{l:PRCRLF}}}\par
\noindent\hyperlink{l:PRDIG}{\texttt{PRDIG}}\hfill\texttt{\$187D~\phantom{000}\llap{\pageref{l:PRDIG}}}\par
\noindent\hyperlink{l:PRDIG1}{\texttt{PRDIG1}}\hfill\texttt{\$1889~\phantom{000}\llap{\pageref{l:PRDIG1}}}\par
\noindent\hyperlink{l:PREC1}{\texttt{PREC1}}\hfill\texttt{\$23B5~\phantom{000}\llap{\pageref{l:PREC1}}}\par
\noindent\hyperlink{l:PREC10}{\texttt{PREC10}}\hfill\texttt{\$23F3~\phantom{000}\llap{\pageref{l:PREC10}}}\par
\noindent\hyperlink{l:PREC2}{\texttt{PREC2}}\hfill\texttt{\$23BC~\phantom{000}\llap{\pageref{l:PREC2}}}\par
\noindent\hyperlink{l:PREC3}{\texttt{PREC3}}\hfill\texttt{\$23C3~\phantom{000}\llap{\pageref{l:PREC3}}}\par
\noindent\hyperlink{l:PREC4}{\texttt{PREC4}}\hfill\texttt{\$23CA~\phantom{000}\llap{\pageref{l:PREC4}}}\par
\noindent\hyperlink{l:PREC5}{\texttt{PREC5}}\hfill\texttt{\$23D1~\phantom{000}\llap{\pageref{l:PREC5}}}\par
\noindent\hyperlink{l:PREC6}{\texttt{PREC6}}\hfill\texttt{\$23D7~\phantom{000}\llap{\pageref{l:PREC6}}}\par
\noindent\hyperlink{l:PREC7}{\texttt{PREC7}}\hfill\texttt{\$23DE~\phantom{000}\llap{\pageref{l:PREC7}}}\par
\noindent\hyperlink{l:PREC8}{\texttt{PREC8}}\hfill\texttt{\$23E5~\phantom{000}\llap{\pageref{l:PREC8}}}\par
\noindent\hyperlink{l:PREC9}{\texttt{PREC9}}\hfill\texttt{\$23EC~\phantom{000}\llap{\pageref{l:PREC9}}}\par
\noindent\hyperlink{l:PRECLV}{\texttt{PRECLV}}\hfill\texttt{\$2705~\phantom{000}\llap{\pageref{l:PRECLV}}}\par
\noindent\hyperlink{l:PRELNK}{\texttt{PRELNK}}\hfill\texttt{\$0E6E~\phantom{000}\llap{\pageref{l:PRELNK}}}\par
\noindent\hyperlink{l:PREPST}{\texttt{PREPST}}\hfill\texttt{\$2077~\phantom{000}\llap{\pageref{l:PREPST}}}\par
\noindent\hyperlink{l:PRIMARY}{\texttt{PRIMARY}}\hfill\texttt{\$257F~\phantom{000}\llap{\pageref{l:PRIMARY}}}\par
\noindent\hyperlink{l:PRINT}{\texttt{PRINT}}\hfill\texttt{\$08D4~\phantom{000}\llap{\pageref{l:PRINT}}}\par
\noindent\hyperlink{l:PRINTS}{\texttt{PRINTS}}\hfill\texttt{\$219F~\phantom{000}\llap{\pageref{l:PRINTS}}}\par
\noindent\hyperlink{l:PRNUM}{\texttt{PRNUM}}\hfill\texttt{\$1862~\phantom{000}\llap{\pageref{l:PRNUM}}}\par
\noindent\hyperlink{l:PRSGN}{\texttt{PRSGN}}\hfill\texttt{\$1860~\phantom{000}\llap{\pageref{l:PRSGN}}}\par
\noindent\hyperlink{l:PRSIN}{\texttt{PRSIN}}\hfill\texttt{\$1855~\phantom{000}\llap{\pageref{l:PRSIN}}}\par
\noindent\hyperlink{l:PRTCOM}{\texttt{PRTCOM}}\hfill\texttt{\$221C~\phantom{000}\llap{\pageref{l:PRTCOM}}}\par
\noindent\hyperlink{l:PRTEXP}{\texttt{PRTEXP}}\hfill\texttt{\$21C0~\phantom{000}\llap{\pageref{l:PRTEXP}}}\par
\noindent\hyperlink{l:PRTFIT}{\texttt{PRTFIT}}\hfill\texttt{\$21F0~\phantom{000}\llap{\pageref{l:PRTFIT}}}\par
\noindent\hyperlink{l:PRTSEP}{\texttt{PRTSEP}}\hfill\texttt{\$223A~\phantom{000}\llap{\pageref{l:PRTSEP}}}\par
\noindent\hyperlink{l:PRTXT}{\texttt{PRTXT}}\hfill\texttt{\$01A8~\phantom{000}\llap{\pageref{l:PRTXT}}}\par
\noindent\hyperlink{l:PUSH0}{\texttt{PUSH0}}\hfill\texttt{\$1C7F~\phantom{000}\llap{\pageref{l:PUSH0}}}\par
\noindent\hyperlink{l:PUTBC}{\texttt{PUTBC}}\hfill\texttt{\$239B~\phantom{000}\llap{\pageref{l:PUTBC}}}\par
\noindent\hyperlink{l:PUTTOK}{\texttt{PUTTOK}}\hfill\texttt{\$2368~\phantom{000}\llap{\pageref{l:PUTTOK}}}\par
\noindent\hyperlink{l:RANDOM}{\texttt{RANDOM}}\hfill\texttt{\$2258~\phantom{000}\llap{\pageref{l:RANDOM}}}\par
\noindent\hyperlink{l:RCVBIT}{\texttt{RCVBIT}}\hfill\texttt{\$061C~\phantom{000}\llap{\pageref{l:RCVBIT}}}\par
\noindent\hyperlink{l:RCVBLK}{\texttt{RCVBLK}}\hfill\texttt{\$05D0~\phantom{000}\llap{\pageref{l:RCVBLK}}}\par
\noindent\hyperlink{l:RCVBYT}{\texttt{RCVBYT}}\hfill\texttt{\$061A~\phantom{000}\llap{\pageref{l:RCVBYT}}}\par
\noindent\hyperlink{l:RCVSTA}{\texttt{RCVSTA}}\hfill\texttt{\$0556~\phantom{000}\llap{\pageref{l:RCVSTA}}}\par
\noindent\hyperlink{l:RDBC}{\texttt{RDBC}}\hfill\texttt{\$1017~\phantom{000}\llap{\pageref{l:RDBC}}}\par
\noindent\hyperlink{l:RDBC1}{\texttt{RDBC1}}\hfill\texttt{\$1016~\phantom{000}\llap{\pageref{l:RDBC1}}}\par
\noindent\hyperlink{l:RDDE}{\texttt{RDDE}}\hfill\texttt{\$101D~\phantom{000}\llap{\pageref{l:RDDE}}}\par
\noindent\hyperlink{l:RDDE1}{\texttt{RDDE1}}\hfill\texttt{\$101C~\phantom{000}\llap{\pageref{l:RDDE1}}}\par
\noindent\hyperlink{l:RDVAL}{\texttt{RDVAL}}\hfill\texttt{\$1EFE~\phantom{000}\llap{\pageref{l:RDVAL}}}\par
\noindent\hyperlink{l:READ}{\texttt{READ}}\hfill\texttt{\$1EC5~\phantom{000}\llap{\pageref{l:READ}}}\par
\noindent\hyperlink{l:READY}{\texttt{READY}}\hfill\texttt{\$00CC~\phantom{000}\llap{\pageref{l:READY}}}\par
\noindent\hyperlink{l:REFADR}{\texttt{REFADR}}\hfill\texttt{\$0D33~\phantom{000}\llap{\pageref{l:REFADR}}}\par
\noindent\hyperlink{l:REM}{\texttt{REM}}\hfill\texttt{\$2276~\phantom{000}\llap{\pageref{l:REM}}}\par
\noindent\hyperlink{l:RESET}{\texttt{RESET}}\hfill\texttt{\$0000~\phantom{000}\llap{\pageref{l:RESET}}}\par
\noindent\hyperlink{l:RESTOR}{\texttt{RESTOR}}\hfill\texttt{\$2286~\phantom{000}\llap{\pageref{l:RESTOR}}}\par
\noindent\hyperlink{l:RETURN}{\texttt{RETURN}}\hfill\texttt{\$22B0~\phantom{000}\llap{\pageref{l:RETURN}}}\par
\noindent\hyperlink{l:RNDADD}{\texttt{RNDADD}}\hfill\texttt{\$32E5~\phantom{000}\llap{\pageref{l:RNDADD}}}\par
\noindent\hyperlink{l:RNDINI}{\texttt{RNDINI}}\hfill\texttt{\$225D~\phantom{000}\llap{\pageref{l:RNDINI}}}\par
\noindent\hyperlink{l:RNDMUL}{\texttt{RNDMUL}}\hfill\texttt{\$32E0~\phantom{000}\llap{\pageref{l:RNDMUL}}}\par
\noindent\hyperlink{l:ROOM}{\texttt{ROOM}}\hfill\texttt{\$339B~\phantom{000}\llap{\pageref{l:ROOM}}}\par
\noindent\hyperlink{l:ROTTXT}{\texttt{ROTTXT}}\hfill\texttt{\$1449~\phantom{000}\llap{\pageref{l:ROTTXT}}}\par
\noindent\hyperlink{l:ROWTAB}{\texttt{ROWTAB}}\hfill\texttt{\$0374~\phantom{000}\llap{\pageref{l:ROWTAB}}}\par
\noindent\hyperlink{l:RST08}{\texttt{RST08}}\hfill\texttt{\$0008~\phantom{000}\llap{\pageref{l:RST08}}}\par
\noindent\hyperlink{l:RST10}{\texttt{RST10}}\hfill\texttt{\$0010~\phantom{000}\llap{\pageref{l:RST10}}}\par
\noindent\hyperlink{l:RST18}{\texttt{RST18}}\hfill\texttt{\$0018~\phantom{000}\llap{\pageref{l:RST18}}}\par
\noindent\hyperlink{l:RST20}{\texttt{RST20}}\hfill\texttt{\$0020~\phantom{000}\llap{\pageref{l:RST20}}}\par
\noindent\hyperlink{l:RST28}{\texttt{RST28}}\hfill\texttt{\$0028~\phantom{000}\llap{\pageref{l:RST28}}}\par
\noindent\hyperlink{l:RST30}{\texttt{RST30}}\hfill\texttt{\$0030~\phantom{000}\llap{\pageref{l:RST30}}}\par
\noindent\hyperlink{l:RST38}{\texttt{RST38}}\hfill\texttt{\$0038~\phantom{000}\llap{\pageref{l:RST38}}}\par
\noindent\hyperlink{l:RSTDAT}{\texttt{RSTDAT}}\hfill\texttt{\$2293~\phantom{000}\llap{\pageref{l:RSTDAT}}}\par
\noindent\hyperlink{l:RUBOUT}{\texttt{RUBOUT}}\hfill\texttt{\$035B~\phantom{000}\llap{\pageref{l:RUBOUT}}}\par
\noindent\hyperlink{l:RUNLIN}{\texttt{RUNLIN}}\hfill\texttt{\$0D71~\phantom{000}\llap{\pageref{l:RUNLIN}}}\par
\noindent\hyperlink{l:RUNSTM}{\texttt{RUNSTM}}\hfill\texttt{\$0D92~\phantom{000}\llap{\pageref{l:RUNSTM}}}\par
\noindent\hyperlink{l:SCRPOS}{\texttt{SCRPOS}}\hfill\texttt{\$0338~\phantom{000}\llap{\pageref{l:SCRPOS}}}\par
\noindent\hyperlink{l:SETDOT}{\texttt{SETDOT}}\hfill\texttt{\$22C7~\phantom{000}\llap{\pageref{l:SETDOT}}}\par
\noindent\hyperlink{l:SETEXT}{\texttt{SETEXT}}\hfill\texttt{\$08BD~\phantom{000}\llap{\pageref{l:SETEXT}}}\par
\noindent\hyperlink{l:SFALSE}{\texttt{SFALSE}}\hfill\texttt{\$2FE9~\phantom{000}\llap{\pageref{l:SFALSE}}}\par
\noindent\hyperlink{l:SINCOS}{\texttt{SINCOS}}\hfill\texttt{\$3DDA~\phantom{000}\llap{\pageref{l:SINCOS}}}\par
\noindent\hyperlink{l:SINSER}{\texttt{SINSER}}\hfill\texttt{\$3FD4~\phantom{000}\llap{\pageref{l:SINSER}}}\par
\noindent\hyperlink{l:SINSGN}{\texttt{SINSGN}}\hfill\texttt{\$3DE0~\phantom{000}\llap{\pageref{l:SINSGN}}}\par
\noindent\hyperlink{l:SKIPS0}{\texttt{SKIPS0}}\hfill\texttt{\$25BF~\phantom{000}\llap{\pageref{l:SKIPS0}}}\par
\noindent\hyperlink{l:SKIPSP}{\texttt{SKIPSP}}\hfill\texttt{\$25BE~\phantom{000}\llap{\pageref{l:SKIPSP}}}\par
\noindent\hyperlink{l:SKPADD}{\texttt{SKPADD}}\hfill\texttt{\$10F9~\phantom{000}\llap{\pageref{l:SKPADD}}}\par
\noindent\hyperlink{l:SKPARG}{\texttt{SKPARG}}\hfill\texttt{\$1B85~\phantom{000}\llap{\pageref{l:SKPARG}}}\par
\noindent\hyperlink{l:SKPELM}{\texttt{SKPELM}}\hfill\texttt{\$0EFB~\phantom{000}\llap{\pageref{l:SKPELM}}}\par
\noindent\hyperlink{l:SKPEXP}{\texttt{SKPEXP}}\hfill\texttt{\$10C3~\phantom{000}\llap{\pageref{l:SKPEXP}}}\par
\noindent\hyperlink{l:SKPFN}{\texttt{SKPFN}}\hfill\texttt{\$10EA~\phantom{000}\llap{\pageref{l:SKPFN}}}\par
\noindent\hyperlink{l:SKPLEN}{\texttt{SKPLEN}}\hfill\texttt{\$10F8~\phantom{000}\llap{\pageref{l:SKPLEN}}}\par
\noindent\hyperlink{l:SKPRET}{\texttt{SKPRET}}\hfill\texttt{\$10E8~\phantom{000}\llap{\pageref{l:SKPRET}}}\par
\noindent\hyperlink{l:SKPSTR}{\texttt{SKPSTR}}\hfill\texttt{\$10F4~\phantom{000}\llap{\pageref{l:SKPSTR}}}\par
\noindent\hyperlink{l:SKPTOK}{\texttt{SKPTOK}}\hfill\texttt{\$10DD~\phantom{000}\llap{\pageref{l:SKPTOK}}}\par
\noindent\hyperlink{l:SMUL}{\texttt{SMUL}}\hfill\texttt{\$2B83~\phantom{000}\llap{\pageref{l:SMUL}}}\par
\noindent\hyperlink{l:SMULDV}{\texttt{SMULDV}}\hfill\texttt{\$2B87~\phantom{000}\llap{\pageref{l:SMULDV}}}\par
\noindent\hyperlink{l:SPACE}{\texttt{SPACE}}\hfill\texttt{\$11EA~\phantom{000}\llap{\pageref{l:SPACE}}}\par
\noindent\hyperlink{l:SPACES}{\texttt{SPACES}}\hfill\texttt{\$224A~\phantom{000}\llap{\pageref{l:SPACES}}}\par
\noindent\hyperlink{l:SQRADD}{\texttt{SQRADD}}\hfill\texttt{\$3E71~\phantom{000}\llap{\pageref{l:SQRADD}}}\par
\noindent\hyperlink{l:SQRBOR}{\texttt{SQRBOR}}\hfill\texttt{\$3E68~\phantom{000}\llap{\pageref{l:SQRBOR}}}\par
\noindent\hyperlink{l:SQRBYT}{\texttt{SQRBYT}}\hfill\texttt{\$3E8D~\phantom{000}\llap{\pageref{l:SQRBYT}}}\par
\noindent\hyperlink{l:SQRCLR}{\texttt{SQRCLR}}\hfill\texttt{\$3E2A~\phantom{000}\llap{\pageref{l:SQRCLR}}}\par
\noindent\hyperlink{l:SQRCY}{\texttt{SQRCY}}\hfill\texttt{\$3E53~\phantom{000}\llap{\pageref{l:SQRCY}}}\par
\noindent\hyperlink{l:SQRDIG}{\texttt{SQRDIG}}\hfill\texttt{\$3E5B~\phantom{000}\llap{\pageref{l:SQRDIG}}}\par
\noindent\hyperlink{l:SQREND}{\texttt{SQREND}}\hfill\texttt{\$3E9E~\phantom{000}\llap{\pageref{l:SQREND}}}\par
\noindent\hyperlink{l:SQRLP}{\texttt{SQRLP}}\hfill\texttt{\$3E37~\phantom{000}\llap{\pageref{l:SQRLP}}}\par
\noindent\hyperlink{l:SQRNXT}{\texttt{SQRNXT}}\hfill\texttt{\$3E86~\phantom{000}\llap{\pageref{l:SQRNXT}}}\par
\noindent\hyperlink{l:SQRRLD}{\texttt{SQRRLD}}\hfill\texttt{\$3E95~\phantom{000}\llap{\pageref{l:SQRRLD}}}\par
\noindent\hyperlink{l:SQRSHL}{\texttt{SQRSHL}}\hfill\texttt{\$3E90~\phantom{000}\llap{\pageref{l:SQRSHL}}}\par
\noindent\hyperlink{l:SQRSUB}{\texttt{SQRSUB}}\hfill\texttt{\$3E40~\phantom{000}\llap{\pageref{l:SQRSUB}}}\par
\noindent\hyperlink{l:STEXE1}{\texttt{STEXE1}}\hfill\texttt{\$0E14~\phantom{000}\llap{\pageref{l:STEXE1}}}\par
\noindent\hyperlink{l:STEXE2}{\texttt{STEXE2}}\hfill\texttt{\$0E42~\phantom{000}\llap{\pageref{l:STEXE2}}}\par
\noindent\hyperlink{l:STKCHK}{\texttt{STKCHK}}\hfill\texttt{\$3389~\phantom{000}\llap{\pageref{l:STKCHK}}}\par
\noindent\hyperlink{l:STKW1}{\texttt{STKW1}}\hfill\texttt{\$1911~\phantom{000}\llap{\pageref{l:STKW1}}}\par
\noindent\hyperlink{l:STKW2}{\texttt{STKW2}}\hfill\texttt{\$199C~\phantom{000}\llap{\pageref{l:STKW2}}}\par
\noindent\hyperlink{l:STMEND}{\texttt{STMEND}}\hfill\texttt{\$218D~\phantom{000}\llap{\pageref{l:STMEND}}}\par
\noindent\hyperlink{l:STOP}{\texttt{STOP}}\hfill\texttt{\$2320~\phantom{000}\llap{\pageref{l:STOP}}}\par
\noindent\hyperlink{l:STRCMP}{\texttt{STRCMP}}\hfill\texttt{\$3012~\phantom{000}\llap{\pageref{l:STRCMP}}}\par
\noindent\hyperlink{l:STRCPN}{\texttt{STRCPN}}\hfill\texttt{\$3068~\phantom{000}\llap{\pageref{l:STRCPN}}}\par
\noindent\hyperlink{l:STRFIL}{\texttt{STRFIL}}\hfill\texttt{\$30FE~\phantom{000}\llap{\pageref{l:STRFIL}}}\par
\noindent\hyperlink{l:STRPSH}{\texttt{STRPSH}}\hfill\texttt{\$2FBB~\phantom{000}\llap{\pageref{l:STRPSH}}}\par
\noindent\hyperlink{l:STRSKP}{\texttt{STRSKP}}\hfill\texttt{\$2F25~\phantom{000}\llap{\pageref{l:STRSKP}}}\par
\noindent\hyperlink{l:STRUE}{\texttt{STRUE}}\hfill\texttt{\$2FF2~\phantom{000}\llap{\pageref{l:STRUE}}}\par
\noindent\hyperlink{l:STTAB1}{\texttt{STTAB1}}\hfill\texttt{\$1A08~\phantom{000}\llap{\pageref{l:STTAB1}}}\par
\noindent\hyperlink{l:STTAB2}{\texttt{STTAB2}}\hfill\texttt{\$1A36~\phantom{000}\llap{\pageref{l:STTAB2}}}\par
\noindent\hyperlink{l:SUBF}{\texttt{SUBF}}\hfill\texttt{\$2ABD~\phantom{000}\llap{\pageref{l:SUBF}}}\par
\noindent\hyperlink{l:SUBI}{\texttt{SUBI}}\hfill\texttt{\$2B62~\phantom{000}\llap{\pageref{l:SUBI}}}\par
\noindent\hyperlink{l:TOFLT}{\texttt{TOFLT}}\hfill\texttt{\$256E~\phantom{000}\llap{\pageref{l:TOFLT}}}\par
\noindent\hyperlink{l:TOINT}{\texttt{TOINT}}\hfill\texttt{\$2561~\phantom{000}\llap{\pageref{l:TOINT}}}\par
\noindent\hyperlink{l:TOKKW}{\texttt{TOKKW}}\hfill\texttt{\$18D4~\phantom{000}\llap{\pageref{l:TOKKW}}}\par
\noindent\hyperlink{l:TOKLEN}{\texttt{TOKLEN}}\hfill\texttt{\$1105~\phantom{000}\llap{\pageref{l:TOKLEN}}}\par
\noindent\hyperlink{l:TOKLIN}{\texttt{TOKLIN}}\hfill\texttt{\$1890~\phantom{000}\llap{\pageref{l:TOKLIN}}}\par
\noindent\hyperlink{l:TOKSTM}{\texttt{TOKSTM}}\hfill\texttt{\$18C7~\phantom{000}\llap{\pageref{l:TOKSTM}}}\par
\noindent\hyperlink{l:TRACE}{\texttt{TRACE}}\hfill\texttt{\$2325~\phantom{000}\llap{\pageref{l:TRACE}}}\par
\noindent\hyperlink{l:TRCPRT}{\texttt{TRCPRT}}\hfill\texttt{\$232F~\phantom{000}\llap{\pageref{l:TRCPRT}}}\par
\noindent\hyperlink{l:TSTBRK}{\texttt{TSTBRK}}\hfill\texttt{\$033E~\phantom{000}\llap{\pageref{l:TSTBRK}}}\par
\noindent\hyperlink{l:TXTBLK}{\texttt{TXTBLK}}\hfill\texttt{\$0995~\phantom{000}\llap{\pageref{l:TXTBLK}}}\par
\noindent\hyperlink{l:TXTEND}{\texttt{TXTEND}}\hfill\texttt{\$0985~\phantom{000}\llap{\pageref{l:TXTEND}}}\par
\noindent\hyperlink{l:TXTINP}{\texttt{TXTINP}}\hfill\texttt{\$09A0~\phantom{000}\llap{\pageref{l:TXTINP}}}\par
\noindent\hyperlink{l:TXTPRT}{\texttt{TXTPRT}}\hfill\texttt{\$0912~\phantom{000}\llap{\pageref{l:TXTPRT}}}\par
\noindent\hyperlink{l:TYPCNV}{\texttt{TYPCNV}}\hfill\texttt{\$2551~\phantom{000}\llap{\pageref{l:TYPCNV}}}\par
\noindent\hyperlink{l:VAR0}{\texttt{VAR0}}\hfill\texttt{\$2CCD~\phantom{000}\llap{\pageref{l:VAR0}}}\par
\noindent\hyperlink{l:VAR1}{\texttt{VAR1}}\hfill\texttt{\$2CE7~\phantom{000}\llap{\pageref{l:VAR1}}}\par
\noindent\hyperlink{l:VAR1F}{\texttt{VAR1F}}\hfill\texttt{\$2CD9~\phantom{000}\llap{\pageref{l:VAR1F}}}\par
\noindent\hyperlink{l:VAR1I}{\texttt{VAR1I}}\hfill\texttt{\$2CDF~\phantom{000}\llap{\pageref{l:VAR1I}}}\par
\noindent\hyperlink{l:VAR1S}{\texttt{VAR1S}}\hfill\texttt{\$2CE4~\phantom{000}\llap{\pageref{l:VAR1S}}}\par
\noindent\hyperlink{l:VAR2}{\texttt{VAR2}}\hfill\texttt{\$2D39~\phantom{000}\llap{\pageref{l:VAR2}}}\par
\noindent\hyperlink{l:VAR2F}{\texttt{VAR2F}}\hfill\texttt{\$2D2B~\phantom{000}\llap{\pageref{l:VAR2F}}}\par
\noindent\hyperlink{l:VAR2I}{\texttt{VAR2I}}\hfill\texttt{\$2D31~\phantom{000}\llap{\pageref{l:VAR2I}}}\par
\noindent\hyperlink{l:VAR2S}{\texttt{VAR2S}}\hfill\texttt{\$2D36~\phantom{000}\llap{\pageref{l:VAR2S}}}\par
\noindent\hyperlink{l:VARIDX}{\texttt{VARIDX}}\hfill\texttt{\$2D11~\phantom{000}\llap{\pageref{l:VARIDX}}}\par
\noindent\hyperlink{l:VARNAM}{\texttt{VARNAM}}\hfill\texttt{\$26CB~\phantom{000}\llap{\pageref{l:VARNAM}}}\par
\noindent\hyperlink{l:VARSIZ}{\texttt{VARSIZ}}\hfill\texttt{\$10B8~\phantom{000}\llap{\pageref{l:VARSIZ}}}\par
\noindent\hyperlink{l:WRBC}{\texttt{WRBC}}\hfill\texttt{\$1029~\phantom{000}\llap{\pageref{l:WRBC}}}\par
\noindent\hyperlink{l:WRBC1}{\texttt{WRBC1}}\hfill\texttt{\$1028~\phantom{000}\llap{\pageref{l:WRBC1}}}\par
\noindent\hyperlink{l:WRDE}{\texttt{WRDE}}\hfill\texttt{\$1023~\phantom{000}\llap{\pageref{l:WRDE}}}\par
\noindent\hyperlink{l:WRDE1}{\texttt{WRDE1}}\hfill\texttt{\$1022~\phantom{000}\llap{\pageref{l:WRDE1}}}\par
\noindent\hyperlink{l:WRDIG}{\texttt{WRDIG}}\hfill\texttt{\$1815~\phantom{000}\llap{\pageref{l:WRDIG}}}\par
\noindent\hyperlink{l:WRZERO}{\texttt{WRZERO}}\hfill\texttt{\$1761~\phantom{000}\llap{\pageref{l:WRZERO}}}\par
\end{multicols}

\end{document}
